%% Custom Pandoc template for PeerJ CS using wlpeerj.cls
%% Usage: pandoc --template=template-peerj.tex paper-peerj.md ...
%%
%% This template maps YAML frontmatter to wlpeerj.cls macros and includes
%% all Pandoc-required LaTeX support (CSL citations, longtable, tight lists,
%% code highlighting, graphics). It absorbs the former preamble.tex formatting.
%%
\PassOptionsToPackage{numbers}{natbib} % citeproc handles citation formatting; prevent natbib author-year conflict
\documentclass[fleqn,10pt,lineno]{wlpeerj}

%% ============================================================================
%% Additional packages for Pandoc output (not included by wlpeerj.cls)
%% ============================================================================

%% Unicode font support via fontspec (XeLaTeX) — fixes Turkish chars and box-drawing glyphs
\usepackage{fontspec}
\setmainfont{Times New Roman}
\setsansfont{Helvetica}
\setmonofont{Menlo}[Scale=0.85]
%% Override wlpeerj.cls author fonts — cls hardcodes \usefont{OT1}{phv}{b}{n} which lacks Turkish chars
\renewcommand\Authfont{\fontsize{12}{14}\selectfont\sffamily\bfseries}
\renewcommand\Affilfont{\fontsize{10}{12}\selectfont\sffamily\bfseries}

%% Tables — Pandoc converts Markdown tables to longtable
\usepackage{longtable}
\usepackage{array}
\newcounter{none} % Pandoc uses \def\LTcaptype{none} for unnumbered tables

%% etoolbox — needed for \AtBeginEnvironment hook
\usepackage{etoolbox}

%% Correct order of tables after \paragraph or \subparagraph
\makeatletter
\patchcmd\longtable{\par}{\if@noskipsec\mbox{}\fi\par}{}{}
\makeatother

%% Allow footnotes in longtable head/foot
\IfFileExists{footnotehyper.sty}{\usepackage{footnotehyper}}{\usepackage{footnote}}
\makesavenoteenv{longtable}

%% Graphics — Pandoc image scaling support
\makeatletter
\newsavebox\pandoc@box
\newcommand*\pandocbounded[1]{% scales image to fit in text height/width
  \sbox\pandoc@box{#1}%
  \Gscale@div\@tempa{\textheight}{\dimexpr\ht\pandoc@box+\dp\pandoc@box\relax}%
  \Gscale@div\@tempb{\linewidth}{\wd\pandoc@box}%
  \ifdim\@tempb\p@<\@tempa\p@\let\@tempa\@tempb\fi%
  \ifdim\@tempa\p@<\p@\scalebox{\@tempa}{\usebox\pandoc@box}%
  \else\usebox{\pandoc@box}%
  \fi%
}
\def\fps@figure{htbp}
\makeatother

%% Tight lists — Pandoc uses this for compact lists
\providecommand{\tightlist}{%
  \setlength{\itemsep}{0pt}\setlength{\parskip}{0pt}}

%% Code highlighting support

%% CSL citation support (for --citeproc)
\NewDocumentCommand\citeproctext{}{}
\NewDocumentCommand\citeproc{mm}{%
  \begingroup\def\citeproctext{#2}\cite{#1}\endgroup}
\makeatletter
 \let\@cite@ofmt\@firstofone
 \def\@biblabel#1{}
 \def\@cite#1#2{{#1\if@tempswa , #2\fi}}
\makeatother
\newlength{\cslhangindent}
\setlength{\cslhangindent}{1.5em}
\newlength{\csllabelwidth}
\setlength{\csllabelwidth}{3em}
\newenvironment{CSLReferences}[2]
 {\begin{list}{}{%
  \setlength{\itemindent}{0pt}
  \setlength{\leftmargin}{0pt}
  \setlength{\parsep}{0pt}
  \ifodd #1
   \setlength{\leftmargin}{\cslhangindent}
   \setlength{\itemindent}{-1\cslhangindent}
  \fi
  \setlength{\itemsep}{#2\baselineskip}}}
 {\end{list}}
\newcommand{\CSLBlock}[1]{\hfill\break\parbox[t]{\linewidth}{\strut\ignorespaces#1\strut}}
\newcommand{\CSLLeftMargin}[1]{\parbox[t]{\csllabelwidth}{\strut#1\strut}}
\newcommand{\CSLRightInline}[1]{\parbox[t]{\linewidth - \csllabelwidth}{\strut#1\strut}}
\newcommand{\CSLIndent}[1]{\hspace{\cslhangindent}#1}

%% ============================================================================
%% Formatting (absorbed from preamble.tex)
%% ============================================================================

%% Reduce font size inside all longtables to prevent overflow
\AtBeginEnvironment{longtable}{%
  \scriptsize%
  \setlength{\tabcolsep}{2pt}%
}

%% Prevent orphan/widow lines (single lines stranded on a page)
\widowpenalties 3 10000 10000 150
\clubpenalties  3 10000 10000 150

%% PeerJ requires left-justified text (not full-width justified)
\raggedright
\setlength{\parindent}{1.5em} % Restore first-line indent zeroed by \raggedright

%% Prevent overfull lines
\setlength{\emergencystretch}{3em}

%% Remove section numbering (Pandoc Markdown uses unnumbered sections)
\setcounter{secnumdepth}{-\maxdimen}

%% ============================================================================
%% Hyperlinks (must be loaded after other packages)
%% ============================================================================
\usepackage{hyperref}
\hypersetup{
  pdftitle={Source Known Identifiers: A Three-Tier Identity System for Distributed Applications},
  pdfauthor={Duran Serkan Kılıç},
  hidelinks,
  pdfcreator={LaTeX via Pandoc}}

%% ============================================================================
%% Extra header includes from YAML (if any)
%% ============================================================================

%% ============================================================================
%% PeerJ metadata from YAML frontmatter
%% ============================================================================
\title{Source Known Identifiers: A Three-Tier Identity System for
Distributed Applications}
\author[1]{Duran Serkan Kılıç}
\affil[1]{Independent Researcher, Ankara, Türkiye}
\corrauthor[1]{Duran Serkan Kılıç}{duran.serkan@outlook.com}

\begin{abstract}
Distributed applications need identifiers that satisfy storage
efficiency, chronological sortability, origin metadata embedding,
zero-lookup verifiability, confidentiality for external consumers, and
multi-century addressability. Based on our literature survey, no
existing scheme provides all six of these identifier properties within a
unified system.

This paper introduces Source Known Identifiers (SKIDs), a three-tier
identity system that projects a single entity identity across trust
boundaries, addressing all six properties. The first tier, Source Known
ID (SKID), is a 64-bit signed integer embedding a timestamp with a
250-millisecond precision, application topology, and a per-entity-type
sequence counter. It serves as the database primary key, providing
compact storage (8 bytes) and natural B-tree ordering for optimized
database indexing. The second tier, Source Known Entity ID (SKEID),
extends the SKID into a 128-bit Universally Unique Identifier (UUID)
compatible value by adding an entity type discriminator, an epoch
selector, and a BLAKE3 keyed message authentication code (MAC). SKEIDs
enable zero-lookup verification of identifier origin, integrity, and
entity type within trusted environments, with a big-endian byte layout
that preserves chronological ordering in lexicographic UUID string
comparisons. The third tier, Secure SKEID, encrypts the entire SKEID
using AES-256 symmetric encryption as a single-block pseudorandom
permutation, producing ciphertext indistinguishable from random bytes
while remaining compatible with standard UUID data-type parsers in
string representation. Deterministic bidirectional transformations
connect all three tiers.

The design employs a defense-in-depth security architecture with
multiple independent verification layers: AES-256 encryption, BLAKE3
keyed MAC, marker byte detection, entity type matching, and record
existence probability. A collision guard mechanism using variant byte
iteration with cryptographic backward verification prevents
misclassification between encrypted and unencrypted forms. An epoch
addressing system spans approximately 17,421 years of total coverage
across 256 configurable epochs.

A reference implementation in C\#/.NET 10 is provided in the open-source
DRN.Framework. BenchmarkDotNet measurements of the reference
implementation (120 iterations, 262,144 invocations per iteration) on
Apple M2 hardware show three performance tiers. Secure SKEID generation
at 544.0 ± 5.7 ns takes approximately 1.4 times as long as UUID Version
7 at 377.5 ± 3.2 ns (a trade-off for providing AES-256 encryption).
SKEID generation at 230.3 ± 3.5 ns is approximately 1.6 times as fast as
UUID Version 7 despite embedding metadata and a BLAKE3 MAC in the same
128-bit footprint. SKID generation at 35.3 ± 1.2 ns is more than 10
times as fast as UUID Version 7 due to deterministic bit-packing without
cryptographic random number generation. All ± values denote 99.9\%
Confidence Interval (CI) error margins.
\end{abstract}

%% ============================================================================
%% Document body
%% ============================================================================
\begin{document}

\flushbottom
\maketitle
\thispagestyle{empty}

\newpage

\section{Introduction}\label{introduction}

\subsection{Problem Statement}\label{problem-statement}

An \emph{entity} is a domain object with a distinct identity that
persists across its lifecycle (Evans, 2003). Its \emph{entity
identifier} is the value that uniquely distinguishes it from all other
entities of the same type. While this terminology is established in
Domain-Driven Design, the underlying requirement of uniquely and
persistently identifying records across distributed systems is universal
regardless of architectural style, design patterns, or definitions.
Specifically, modern distributed applications require entity identifiers
that serve multiple roles simultaneously such as database primary keys,
inter-application correlation tokens, and externally visible resource
handles. However, existing identifier schemes force a choice between
conflicting properties.

Database sequences offer compact storage (4 or 8 bytes) and natural
ordering but expose creation patterns and cannot be safely passed to
external consumers. UUID Version 4 (Davis, Peabody \& Leach, 2024)
provides 122 bits of randomness and global uniqueness but sacrifices
ordering. UUID Version 7 (Davis, Peabody \& Leach, 2024) restores
time-ordering with a 48-bit Unix timestamp but exposes creation
patterns. Snowflake-style identifiers (Twitter Engineering, 2010) embed
timestamp and worker topology in 64 bits but again expose creation
patterns and lose multi-century addressability.

None of these schemes offer a mechanism for zero-lookup verification.
Chronological sortability and confidentiality are also inherently
contradictory properties, and dual-identifier alternatives (e.g.,
integer primary key paired with UUID external-facing identifier) lose
storage efficiency by maintaining two separate identifier columns and
their associated indexes.

A multi-tier identifier system can address these challenges by providing
different identifier properties for different trust levels. The proposed
solution defines the following identifier properties as desired and
achievable in a multi-tier distributed identity system.

\begin{enumerate}
\def\labelenumi{\arabic{enumi}.}
\tightlist
\item
  Storage efficiency
\item
  Chronological sortability
\item
  Origin metadata embedding
\item
  Zero-lookup verifiability
\item
  Confidentiality for external consumers
\item
  Multi-century addressability
\end{enumerate}

Table 1 summarizes identifier property satisfaction across existing
schemes and the SKID system.

\textbf{Table 1:} Identifier property satisfaction across schemes.

{\def\LTcaptype{none} % do not increment counter
\begin{longtable}[]{@{}
  >{\raggedright\arraybackslash}p{(\linewidth - 16\tabcolsep) * \real{0.2520}}
  >{\raggedright\arraybackslash}p{(\linewidth - 16\tabcolsep) * \real{0.1057}}
  >{\raggedright\arraybackslash}p{(\linewidth - 16\tabcolsep) * \real{0.0894}}
  >{\raggedright\arraybackslash}p{(\linewidth - 16\tabcolsep) * \real{0.0894}}
  >{\raggedright\arraybackslash}p{(\linewidth - 16\tabcolsep) * \real{0.0894}}
  >{\raggedright\arraybackslash}p{(\linewidth - 16\tabcolsep) * \real{0.0894}}
  >{\raggedright\arraybackslash}p{(\linewidth - 16\tabcolsep) * \real{0.0813}}
  >{\raggedright\arraybackslash}p{(\linewidth - 16\tabcolsep) * \real{0.0894}}
  >{\raggedright\arraybackslash}p{(\linewidth - 16\tabcolsep) * \real{0.1138}}@{}}
\toprule\noalign{}
\begin{minipage}[b]{\linewidth}\raggedright
Property
\end{minipage} & \begin{minipage}[b]{\linewidth}\raggedright
SKID System
\end{minipage} & \begin{minipage}[b]{\linewidth}\raggedright
UUID V4
\end{minipage} & \begin{minipage}[b]{\linewidth}\raggedright
UUID V7
\end{minipage} & \begin{minipage}[b]{\linewidth}\raggedright
Snowflake
\end{minipage} & \begin{minipage}[b]{\linewidth}\raggedright
ULID
\end{minipage} & \begin{minipage}[b]{\linewidth}\raggedright
CUID2
\end{minipage} & \begin{minipage}[b]{\linewidth}\raggedright
KSUID
\end{minipage} & \begin{minipage}[b]{\linewidth}\raggedright
DB Sequence
\end{minipage} \\
\midrule\noalign{}
\endhead
\bottomrule\noalign{}
\endlastfoot
Storage efficiency & 8/16 B & 16 B & 16 B & 8 B & 16 B & 24 chars & 20 B
& 4/8 B \\
Chronological sortability & Yes & No & Yes & Yes & Yes & No & Yes &
Yes \\
Origin metadata embedding & Yes & No & No & Partial & No & No & No &
No \\
Zero-lookup verifiability & Yes (MAC) & No & No & No & No & No & No &
No \\
Confidentiality (external) & Encryption & Randomness & No & No & No &
Hashing & No & No \\
Multi-century addressability & Yes (SKEID) & N/A & Yes & No & Yes & N/A
& Partial & N/A \\
\end{longtable}
}

Table 2 provides a detailed feature comparison of these and additional
identifier schemes.

\newpage

\subsection{Motivation}\label{motivation}

The necessity for a unified, multi-tier identifier system arises from
the conflicting architectural requirements of modern distributed
applications. Data in these systems continuously flows across
persistence, internal, and external trust tiers. Each tier imposes
different constraints on identifier design. Existing schemes and
dual-identifier patterns force systems to compromise at least one of
these requirements.

At the persistence tier, B-tree index performance demands sequential or
time-ordered primary keys (Davis, Peabody \& Leach, 2024, sec. 2.1 and
6.13). However, exposing these identifiers through external APIs creates
an attack surface for Insecure Direct Object Reference (IDOR)
vulnerabilities (MITRE Corporation, 2025a). Exposure also enables
adversarial inference of generation velocity and record volume (the
German Tank Problem (Ruggles \& Brodie, 1947)). As a workaround, the
dual-identifier pattern with an integer primary key alongside a random
UUID alternate key compromises storage efficiency.

The full cost of such workarounds extends beyond the identifier columns
themselves. For instance, a conventional schema using an auto-increment
primary key (8 bytes), a UUID external identifier (16 bytes), and a
\texttt{created\_at} timestamp (8 bytes) requires 32 bytes of column
data per record. In a fully indexed schema, each column demands a
separate B-tree index, resulting in three indexes and their associated
maintenance operations such as vacuum, reindex, and statistics
collection. An identifier that embeds a timestamp and derives a
UUID-compatible external representation deterministically at application
runtime consolidates all three concerns into a single 8-byte primary key
column with a single index. Under this fully indexed baseline, this is a
75\% reduction in per-record column overhead and a reduction from three
single-column indexes to one. This calculation accounts only for the
single-column indexes on each field. Additional composite indexes
involving these fields would amplify the savings further.

A downstream application receiving an identifier must verify its
authenticity (MITRE Corporation, 2025b). Without embedded verification
metadata, this validation requires a query or cache lookup per
identifier. An identifier carrying a MAC would eliminate I/O-bound
validation overhead and enable zero-lookup verification.

Decentralized identifier generation across globally distributed networks
demands chronologically sortable (Davis, Peabody \& Leach, 2024)
identifiers. When data from independent generators must be merged years
or decades after creation, the identifier scheme must guarantee
sort-order consistency beyond the operational lifespan of individual
system components. This is especially critical for any system requiring
long-term data retention and historical analysis.

These conflicting trust-boundary constraints motivate a multi-tier
identifier system where a single entity identity is projected to satisfy
each trust boundary.

\subsection{Contributions}\label{contributions}

This paper makes the following contributions.

\begin{enumerate}
\def\labelenumi{\arabic{enumi}.}
\tightlist
\item
  \textbf{Design of a three-tier identity system} that satisfies all six
  desired identifier properties through a layered architecture aligned
  with trust boundaries (database, trusted internal, and external).
\item
  \textbf{A defense-in-depth security architecture} combining AES-256
  encryption, BLAKE3 keyed MAC, marker byte detection, entity type
  matching, and a novel collision guard mechanism with cryptographic
  backward verification.
\item
  \textbf{An epoch addressing system} spanning approximately 17,421
  years of total coverage through 256 configurable epochs of \(2^{31}\)
  seconds each.
\item
  \textbf{A reference implementation} in C\#/.NET 10 with open-source
  code, integration and unit test coverage, and BenchmarkDotNet
  performance data demonstrating competitive or superior performance
  compared to standard UUID generation.
\end{enumerate}

\subsection{Related Work}\label{related-work}

Distributed identifier schemes have evolved considerably since the
original UUID specification. RFC 9562 (Davis, Peabody \& Leach, 2024)
standardized UUID versions 1, 3--8, with Version 7 introducing
timestamp-ordered UUIDs that address the B-tree fragmentation problems
of random Version 4 UUIDs. UUID Version 7 uses a 48-bit Unix timestamp
prefix followed by random bits, providing millisecond-precision ordering
and global uniqueness.

Twitter's Snowflake architecture (Twitter Engineering, 2010) pioneered
timestamp-prefixed, worker-partitioned 64-bit identifiers for
high-throughput systems. Snowflake identifiers embed a 41-bit timestamp
with millisecond precision, a 10-bit machine identifier, and a 12-bit
sequence number, enabling approximately 4,096 identifiers per
millisecond per machine. Its custom epoch starts from November 4, 2010
and extends to July 10, 2080 (\textasciitilde69.7 years).

ULIDs (Feerasta, 2016) provide a 128-bit value composed of a 48-bit
timestamp with millisecond precision. Its Unix epoch starts from January
1, 1970 and extends to \textasciitilde10889 CE (\textasciitilde8,919
years). ULIDs are encoded as 26-character Crockford Base32 strings that
sort lexicographically by creation time, targeting environments where
string-based identifiers are standard.

The original CUID (Elliott, 2012) was a collision-resistant identifier
specification that used a k-sortable, timestamp-prefixed structure. It
was deprecated by its author due to security concerns. The CUID
deprecation notice warns that ``all monotonically increasing
(auto-increment, k-sortable), and timestamp-based ids share the security
issues with Cuid'' and further states that ``UUID V6-V8 are also
insecure because they leak information which could be used to exploit
systems or violate user privacy'' (Elliott, 2012).

CUID2 (Elliott, 2022), the successor to CUID, takes a fundamentally
different approach. CUID2 intentionally removed timestamps from the
identifier for security reasons and instead recommends a separate
\texttt{createdAt} column for time-based sorting, adding per-record
storage overhead that timestamp-embedding schemes avoid. CUID2 generates
identifiers by using independent entropy sources then hashing the
concatenation with SHA3. This produces identifiers with strong collision
resistance but relies on probabilistic uniqueness rather than
deterministic construction.

KSUIDs (Segment, 2017) (K-Sortable Unique Identifier) provides a 160-bit
(20-byte) value composed of a 32-bit timestamp. It has 1-second
precision, a custom epoch from May 13, 2014 to June 19, 2150, and a
128-bit cryptographically random payload. KSUIDs are encoded as
27-character Base62 strings that sort lexicographically by creation
time. The 128-bit random payload provides stronger collision resistance
than UUID V4's 122 random bits. The string-first design targets
application-layer identifiers rather than database primary keys where
compact binary representation is critical for B-tree performance.

None of these approaches offers integrity verification. Table 2 presents
a detailed feature comparison.

\textbf{Table 2:} Feature comparison of distributed identifier schemes.

{\def\LTcaptype{none} % do not increment counter
\begin{longtable}[]{@{}
  >{\raggedright\arraybackslash}p{(\linewidth - 16\tabcolsep) * \real{0.2114}}
  >{\raggedright\arraybackslash}p{(\linewidth - 16\tabcolsep) * \real{0.1789}}
  >{\raggedright\arraybackslash}p{(\linewidth - 16\tabcolsep) * \real{0.0976}}
  >{\raggedright\arraybackslash}p{(\linewidth - 16\tabcolsep) * \real{0.0732}}
  >{\raggedright\arraybackslash}p{(\linewidth - 16\tabcolsep) * \real{0.0894}}
  >{\raggedright\arraybackslash}p{(\linewidth - 16\tabcolsep) * \real{0.0813}}
  >{\raggedright\arraybackslash}p{(\linewidth - 16\tabcolsep) * \real{0.0813}}
  >{\raggedright\arraybackslash}p{(\linewidth - 16\tabcolsep) * \real{0.0813}}
  >{\raggedright\arraybackslash}p{(\linewidth - 16\tabcolsep) * \real{0.1057}}@{}}
\toprule\noalign{}
\begin{minipage}[b]{\linewidth}\raggedright
Feature
\end{minipage} & \begin{minipage}[b]{\linewidth}\raggedright
SKID System
\end{minipage} & \begin{minipage}[b]{\linewidth}\raggedright
UUID V4
\end{minipage} & \begin{minipage}[b]{\linewidth}\raggedright
UUID V7
\end{minipage} & \begin{minipage}[b]{\linewidth}\raggedright
Snowflake
\end{minipage} & \begin{minipage}[b]{\linewidth}\raggedright
ULID
\end{minipage} & \begin{minipage}[b]{\linewidth}\raggedright
CUID2
\end{minipage} & \begin{minipage}[b]{\linewidth}\raggedright
KSUID
\end{minipage} & \begin{minipage}[b]{\linewidth}\raggedright
DB Sequence
\end{minipage} \\
\midrule\noalign{}
\endhead
\bottomrule\noalign{}
\endlastfoot
Chronological ordering & 250ms & No & ms & ms & ms & No & second &
Numeric \\
Embedded timestamp & Yes & No & Yes & Yes & Yes & No & Yes & No \\
Multi-century addressing & Yes (SKEID) & N/A & Yes & No & Yes & N/A &
Partial & N/A \\
Entity type & Yes (8-bit) & No & No & No & No & No & No & No \\
Integrity check & Yes (BLAKE3 MAC) & No & No & No & No & No & No & No \\
Confidentiality & Encryption (AES-256) & Randomness & No & No & No &
Hashing & No & No \\
Zero-lookup validation & Yes & No & No & No & No & No & No & No \\
DB primary key size & 8 B (SKID) & 16 B & 16 B & 8 B & 16 B & 24 chars &
20 B & 4/8 B \\
External ID size & 16 B (SKEID) & 16 B & 16 B & 8 B & 26 chars & 24
chars & 27 chars & 4/8 B \\
UUID compatible & Yes (SKEID, UUID V8) & Yes & Yes & No & No & No & No &
No \\
Key rotation & Yes (key-ring) & N/A & N/A & N/A & N/A & N/A & N/A &
N/A \\
\end{longtable}
}

\section{Methods}\label{methods}

\subsection{Three-Tier Identity Model}\label{three-tier-identity-model}

SKIDs address the identifier gap through a three-tier identity model
aligned with three trust boundaries commonly found in distributed
architectures.

\begin{enumerate}
\def\labelenumi{\arabic{enumi}.}
\item
  \textbf{Database Tier:} A 64-bit Source Known ID (SKID) serves as the
  primary key. It embeds a 250-millisecond-precision timestamp,
  application topology (application identifier, instance identifier),
  and a per-entity-type sequence counter, providing natural ordering and
  compact B-tree indexing.
\item
  \textbf{Trusted Environment Tier:} A 128-bit Source Known Entity ID
  (SKEID) encodes the SKID, the entity type, epoch metadata, a keyed MAC
  for integrity, and RFC 9562 Version 8 compatible markers. Applications
  can validate an identifier's origin, entity type, and integrity
  without a database lookup.
\item
  \textbf{External Tier:} A 128-bit Secure Source Known Entity ID
  (Secure SKEID) encrypts the entire SKEID block using AES-256,
  preventing information leakage to untrusted consumers while remaining
  compatible with standard UUID data-type parsers in string
  representation (see Limitation 6).
\end{enumerate}

The three tiers are not three separate identifiers but one entity
identity projected for different trust boundaries. The database stores
the compact SKID, trusted applications derive the SKEID on demand, and
external consumers receive the Secure SKEID. Deterministic bidirectional
transformations between tiers allow any representation to be converted
to any other, given the appropriate cryptographic keys. These
transformations are sub-microsecond CPU operations (217--544 ns, see
Performance Results) requiring no I/O or external lookup. Figure 1
illustrates the data transformations between the three tiers.

\begin{verbatim}
   Database              Trusted Internal          External / Public
  ┌──────────┐          ┌───────────────┐          ┌────────────────┐
  │ SKID     │ Generate │ SKEID         │ ToSecure │ Secure SKEID   │
  │ Sortable │─────────▶│ BLAKE3 MAC    │─────────▶│ AES Encrypted  │
  │ 64-bit   │◀─────────│ 128-bit       │◀─────────│ 128-bit        │
  └──────────┘  Parse   └───────────────┘ ToPlain  └────────────────┘
\end{verbatim}

\textbf{Figure 1:} Three-tier identity model: bidirectional data
transformations.

\subsection{Source Known ID (SKID) 64-bit
Specification}\label{source-known-id-skid-64-bit-specification}

\subsubsection{Bit Layout}\label{bit-layout}

A SKID is a 64-bit signed integer with the following field layout
(Figure 2), packed most-significant bit first.

\begin{verbatim}
 0                   1                   2                   3
 0 1 2 3 4 5 6 7 8 9 0 1 2 3 4 5 6 7 8 9 0 1 2 3 4 5 6 7 8 9 0 1
+-+-+-+-+-+-+-+-+-+-+-+-+-+-+-+-+-+-+-+-+-+-+-+-+-+-+-+-+-+-+-+-+
|S|                    Timestamp (T - 32 bits)                   
+-+-+-+-+-+-+-+-+-+-+-+-+-+-+-+-+-+-+-+-+-+-+-+-+-+-+-+-+-+-+-+-+
 T| App ID (7)  | Inst (6)  |        SequenceId (18 bits)       |
+-+-+-+-+-+-+-+-+-+-+-+-+-+-+-+-+-+-+-+-+-+-+-+-+-+-+-+-+-+-+-+-+
\end{verbatim}

\textbf{Figure 2:} SKID 64-bit field layout (IETF RFC diagram
convention).

Total: 1 (Sign) + 32 (Timestamp) + 7 (App ID) + 6 (App Instance Id) + 18
(Sequence Id) = 64 bits.

\newpage

Table 3 details the corresponding fields.

\textbf{Table 3:} SKID field definitions.

{\def\LTcaptype{none} % do not increment counter
\begin{longtable}[]{@{}
  >{\raggedright\arraybackslash}p{(\linewidth - 6\tabcolsep) * \real{0.1597}}
  >{\raggedright\arraybackslash}p{(\linewidth - 6\tabcolsep) * \real{0.0672}}
  >{\raggedright\arraybackslash}p{(\linewidth - 6\tabcolsep) * \real{0.0756}}
  >{\raggedright\arraybackslash}p{(\linewidth - 6\tabcolsep) * \real{0.6975}}@{}}
\toprule\noalign{}
\begin{minipage}[b]{\linewidth}\raggedright
Field
\end{minipage} & \begin{minipage}[b]{\linewidth}\raggedright
Bit(s)
\end{minipage} & \begin{minipage}[b]{\linewidth}\raggedright
Width
\end{minipage} & \begin{minipage}[b]{\linewidth}\raggedright
Purpose
\end{minipage} \\
\midrule\noalign{}
\endhead
\bottomrule\noalign{}
\endlastfoot
Sign / Epoch half & 63 & 1 bit & Epoch half indicator for smooth epoch
transitions \\
Timestamp & 62--31 & 32 bits & 250ms tick precision, epoch-relative
(approximately 68 years per epoch) \\
App Id & 30--24 & 7 bits & Application identifier (maximum 127) \\
App Instance Id & 23--18 & 6 bits & Instance discriminator (maximum
63) \\
Sequence & 17--0 & 18 bits & Per-instance monotonic counter (262,144 per
250ms tick; 1,048,576 per second) \\
\end{longtable}
}

The sign bit controls epoch half selection. When set to 1, the resulting
value is negative, covering the first approximately 34 years of the
epoch. When set to 0, the value becomes positive, extending coverage for
the second approximately 34 years. Together, the two halves span
approximately 68 years per epoch while preserving monotonic sort order
in signed 64-bit representation, since negative values sort before
positive values. This timestamp-leading layout produces a
chronologically sortable SKID optimized for internal database indexing
and ordering.

The 32-bit timestamp field stores unsigned 250-millisecond ticks elapsed
since the configured epoch, giving four ticks per second. Sub-second
ordering at 250ms granularity eliminates coarse-grained temporal
ambiguity while preserving throughput within a compact 64-bit layout.
The 18-bit sequence counter allows 262,144 unique identifiers per 250ms
tick (1,048,576 per second) per instance. Implementations must reset the
sequence when the timestamp advances. If the sequence is exhausted
within a single tick, the generator must wait for the next tick before
issuing new identifiers, applying backpressure to maintain uniqueness
guarantees.

\subsubsection{Generation Algorithm}\label{generation-algorithm}

The SKID generation procedure operates as follows. Given
\texttt{entityType}, \texttt{appId}, \texttt{appInstanceId}, and a
configured \texttt{epoch}:

\begin{enumerate}
\def\labelenumi{\arabic{enumi}.}
\tightlist
\item
  Compute \texttt{elapsedTicks} as the floor of 250-millisecond ticks
  elapsed since the epoch (four ticks per second).
\item
  Assert that the elapsed ticks are within the epoch range (0 to
  \(2^{33} - 1\) ticks, covering both epoch halves).
\item
  Determine the epoch half:

  \begin{enumerate}
  \def\labelenumii{\arabic{enumii}.}
  \tightlist
  \item
    If \texttt{elapsedTicks} \(< 2^{32}\), the current half is the first
    half (sign bit = 1, producing a negative SKID). Otherwise, the
    current half is the second half (sign bit = 0, producing a positive
    SKID).
  \item
    Compute \texttt{timestamp} = \texttt{elapsedTicks} mod \(2^{32}\)
    (equivalently, \texttt{elapsedTicks} AND \(2^{32} - 1\)).
  \end{enumerate}
\item
  Obtain the next sequence value from the per-entity-type atomic counter
  (resetting on timestamp change).
\item
  Pack fields into a 64-bit signed integer: sign bit at position 63,
  timestamp at positions 62--31, App ID at positions 30--24, App
  Instance ID at positions 23--18, and Sequence ID at positions 17--0.
\item
  Return the packed value.
\end{enumerate}

\subsubsection{Clock Drift Protection}\label{clock-drift-protection}

The system handles backward time jumps at two levels. Minor drifts
freeze the timestamp until the wall clock catches up, allowing the
sequence counter to continue advancing within the frozen tick. The
reference implementation uses 5 seconds as freeze threshold. Drifts
beyond freeze threshold are considered as critical and force the
application instance to shut down and restart with a new instance ID.
This dual-threshold mechanism prevents duplicate or out-of-order
identifiers without requiring external coordination.

\newpage

\subsection{Source Known Entity ID (SKEID) 128-bit
Specification}\label{source-known-entity-id-skeid-128-bit-specification}

\subsubsection{Byte Layout}\label{byte-layout}

An SKEID occupies 16 bytes (128 bits), structured as shown in Figure 3.

\begin{verbatim}
Byte:  0  1  2  3  4  5  6  7  8  9 10 11 12 13 14 15
     +--+--+--+--+--+--+--+--+--+--+--+--+--+--+--+--+
     |EP| SKID Upper|S0|VE|ET|VA|S1|S2|S3|    MAC    |
     +--+--+--+--+--+--+--+--+--+--+--+--+--+--+--+--+
\end{verbatim}

\textbf{Figure 3:} SKEID 16-byte field layout (big-endian, RFC 9562
network byte order).

Table 4 details the byte-level layout of the SKEID fields.

\textbf{Table 4:} SKEID byte layout.

{\def\LTcaptype{none} % do not increment counter
\begin{longtable}[]{@{}
  >{\raggedright\arraybackslash}p{(\linewidth - 6\tabcolsep) * \real{0.1889}}
  >{\raggedright\arraybackslash}p{(\linewidth - 6\tabcolsep) * \real{0.1000}}
  >{\raggedright\arraybackslash}p{(\linewidth - 6\tabcolsep) * \real{0.1000}}
  >{\raggedright\arraybackslash}p{(\linewidth - 6\tabcolsep) * \real{0.6111}}@{}}
\toprule\noalign{}
\begin{minipage}[b]{\linewidth}\raggedright
Field
\end{minipage} & \begin{minipage}[b]{\linewidth}\raggedright
Byte(s)
\end{minipage} & \begin{minipage}[b]{\linewidth}\raggedright
Width
\end{minipage} & \begin{minipage}[b]{\linewidth}\raggedright
Purpose
\end{minipage} \\
\midrule\noalign{}
\endhead
\bottomrule\noalign{}
\endlastfoot
Epoch & 0 & 8 bits & Epoch index (highest lexicographic priority) \\
SKID upper half & 1--4 & 32 bits & SKID bits 63--32, sign-toggled for
lexicographic order \\
SKID low byte 0 & 5 & 8 bits & MSB of SKID lower half (timestamp LSB,
App ID) \\
Version marker & 6 & 8 bits & \texttt{0x8D} (UUID V8, RFC 9562 Section
5.8 octet 6) \\
Entity Type & 7 & 8 bits & Domain entity classification \\
Variant marker & 8 & 8 bits & \texttt{0x8D} (RFC 9562 Section 4.1 octet
8 variant) \\
SKID low bytes & 9--11 & 24 bits & Remaining SKID lower half \\
MAC & 12--15 & 32 bits & BLAKE3 keyed MAC (contiguous) \\
\end{longtable}
}

All fields are encoded in big-endian (network byte order) per RFC 9562.
The epoch byte occupies byte 0 to ensure that higher epoch values sort
lexicographically after lower epoch values, regardless of timestamp. The
upper SKID half at bytes 1--4 is sign-toggled (XOR with
\texttt{0x80000000}) before encoding. This converts the signed SKID sort
order (negative before positive) to unsigned byte sort order, preserving
chronological ordering in lexicographic comparisons across both epoch
halves. The most significant byte of the lower SKID half occupies byte
5, immediately after the upper half, ensuring that the timestamp
least-significant bit and leading application topology bits participate
in lexicographic comparison before the entity type at byte 7. The
remaining lower SKID half is split around the variant marker at byte 8,
with bytes 9--11 holding the last three bytes. Since the version marker,
entity type, and variant marker are constant for all plain SKEIDs of a
given type, lexicographic comparison of the split lower half operates
correctly.

The marker bytes at positions 6 and 8 serve dual purposes.

\begin{enumerate}
\def\labelenumi{\arabic{enumi}.}
\tightlist
\item
  \textbf{UUID V8 compatibility} makes the SKEID a valid RFC 9562 UUID.
\item
  \textbf{Detection signal} enables the parser to distinguish plaintext
  SKEIDs from encrypted Secure SKEIDs.
\end{enumerate}

\subsubsection{MAC Computation}\label{mac-computation}

Integrity verification relies on BLAKE3 keyed MAC (O'Connor et al.,
2020). The computation procedure clears the four MAC bytes (positions
12--15) to zero, then computes a 4-byte BLAKE3 keyed MAC over the full
16-byte buffer. The resulting 4 bytes are written into the contiguous
MAC positions (bytes 12--15).

BLAKE3 was selected for performance and security margin. Our
BenchmarkDotNet measurements show BLAKE3 keyed hashing is approximately
3.5 times faster than HMAC-SHA-256 on small inputs (see supplementary
file \texttt{supplementary-hash-benchmark-report.html}).

Generated MAC is truncated to 32 bits, offering a \(1/2^{32}\)
false-positive rate. This truncation is acceptable within the
defense-in-depth architecture where the MAC is one of multiple
verification layers, not the sole security mechanism.

The epoch byte (byte 0) is not cleared before MAC computation. It
participates in the MAC input, adding tamper-resistance for the epoch
field.

\newpage

\subsubsection{Generation Algorithm}\label{generation-algorithm-1}

Given a 64-bit SKID, epoch, entity type, and MAC key, the SKEID
generation procedure:

\begin{enumerate}
\def\labelenumi{\arabic{enumi}.}
\tightlist
\item
  Allocate a 16-byte buffer initialized to zero.
\item
  Split the SKID into upper and lower 32-bit halves. Toggle the most
  significant bit of the upper half (XOR with \texttt{0x80000000}) and
  write it in big-endian to bytes 1--4.
\item
  Write the most significant byte of the SKID lower half to byte 5.
\item
  Write the epoch byte to byte 0, the version marker \texttt{0x8D} to
  byte 6 (RFC 9562 Section 5.8 octet 6), entity type to byte 7, and
  variant marker to byte 8 (RFC 9562 Section 4.1 octet 8).
\item
  Write the remaining three bytes of the SKID lower half to bytes 9--11.
\item
  Compute the BLAKE3 keyed MAC and write it to bytes 12--15.
\item
  Construct the UUID from the 16-byte buffer, interpreting all fields in
  big-endian (network byte order) per RFC 9562, and return the result.
\end{enumerate}

\subsubsection{Auto-Detection Parse
Logic}\label{auto-detection-parse-logic}

The parse algorithm automatically determines whether a UUID is a
plaintext SKEID, a Secure SKEID, or an unrecognized value.

\begin{enumerate}
\def\labelenumi{\arabic{enumi}.}
\tightlist
\item
  Convert the UUID to a 16-byte big-endian byte array (network byte
  order, per RFC 9562).
\item
  Check for exact primary marker bytes (\texttt{0x8D} at position 6,
  \texttt{0x8D} at position 8) in the byte array.
\item
  If both exact markers are present: attempt plaintext verification. If
  MAC verification succeeds, return as plaintext SKEID.
\item
  If exact markers are absent or plaintext verification fails: decrypt
  the 16-byte block with the AES key.
\item
  Check the decrypted plaintext for markers using a wider acceptance
  rule: \texttt{guidBytes{[}6{]}\ ==\ 0x8D} (exact version match) AND
  \texttt{(guidBytes{[}8{]}\ \&\ 0xC0)\ ==\ 0x80} (RFC 9562 §4.1 variant
  range \texttt{0x80}--\texttt{0xBF}). The wider variant acceptance is
  necessary because the collision guard may have incremented the variant
  byte beyond \texttt{0x8D} during generation.
\item
  If the recovered variant byte is within the RFC 9562 variant range but
  less than \texttt{0x8D}, return INVALID. The generator never produces
  variant bytes below \texttt{0x8D}; such a value indicates tampering or
  an unrelated UUID.
\item
  If the recovered variant byte exceeds \texttt{0x8D}, perform backward
  collision-guard verification as specified in the Backward Verification
  Algorithm section. If verification fails, return INVALID.
\item
  If MAC verifies the decrypted plaintext, return the parsed SKEID with
  a secure-origin flag.
\item
  Otherwise, return INVALID.
\end{enumerate}

The probability of a random or encrypted byte sequence coincidentally
containing the marker bytes \texttt{0x8D8D} at the exact positions is
approximately \(1/65{,}536\). When this occurs, the algorithm gracefully
falls back to the decryption path with no data corruption.

\subsection{Secure SKEID 128-bit
Specification}\label{secure-skeid-128-bit-specification}

\subsubsection{AES-256-ECB Single-Block
Encryption}\label{aes-256-ecb-single-block-encryption}

Table 5 shows the Secure SKEID byte layout after AES-256 encryption.

\textbf{Table 5:} Secure SKEID byte layout.

{\def\LTcaptype{none} % do not increment counter
\begin{longtable}[]{@{}
  >{\raggedright\arraybackslash}p{(\linewidth - 6\tabcolsep) * \real{0.1889}}
  >{\raggedright\arraybackslash}p{(\linewidth - 6\tabcolsep) * \real{0.1000}}
  >{\raggedright\arraybackslash}p{(\linewidth - 6\tabcolsep) * \real{0.1111}}
  >{\raggedright\arraybackslash}p{(\linewidth - 6\tabcolsep) * \real{0.6000}}@{}}
\toprule\noalign{}
\begin{minipage}[b]{\linewidth}\raggedright
Field
\end{minipage} & \begin{minipage}[b]{\linewidth}\raggedright
Byte(s)
\end{minipage} & \begin{minipage}[b]{\linewidth}\raggedright
Width
\end{minipage} & \begin{minipage}[b]{\linewidth}\raggedright
Purpose
\end{minipage} \\
\midrule\noalign{}
\endhead
\bottomrule\noalign{}
\endlastfoot
Ciphertext & 0--15 & 128 bits & AES-256 encrypted SKEID (pseudorandom
bytes) \\
\end{longtable}
}

As a single AES block, the entire SKEID is encrypted. No internal
structure is visible to external consumers.

A Secure SKEID is produced by encrypting the entire 16-byte SKEID
plaintext using AES-256 (National Institute of Standards and Technology,
2001) as a single-block cipher. Under the standard PRP assumption
(Mouha, 2021), AES-256 applied to a single 128-bit block functions as a
pseudorandom permutation (PRP), meaning every distinct plaintext maps to
a distinct ciphertext and vice versa. A pseudorandom function (PRF)
similarly produces outputs indistinguishable from random. The PRP/PRF
switching lemma establishes a \(2^{n/2}\) distinguishing bound for
\(n\)-bit block ciphers (Bellare, Krovetz \& Rogaway, 1998). For AES
with \(n = 128\), a single-block PRP encryption is therefore
indistinguishable from a random function up to approximately \(2^{64}\)
queries. This is beyond any practical identifier generation volume.

SKEID encryption always operates on exactly one block, the multi-block
weakness of ECB mode does not apply. For a single block, ECB is
mathematically identical to Cipher Block Chaining (CBC) with a zero
initialization vector:
\(C = \text{AES}(Key, P \oplus 0) = \text{AES}(Key, P)\). No nonce is
required, and there is no nonce-reuse vulnerability. ECB avoids the
allocation and XOR overhead of a CBC initialization vector that would
provide no additional security for single-block operations.

Quantum attacks on AES via Grover's algorithm reduce the effective
security of AES-256 to 128 bits, which still remains computationally
infeasible (Bonnetain, Naya-Plasencia \& Schrottenloher, 2019).

\subsubsection{Key Separation}\label{key-separation}

The SKEID system uses two cryptographically independent keys derived
from each key-ring entry. A MAC key is used for BLAKE3 keyed MAC
(integrity verification) and an AES key is used for AES-256-ECB
(confidentiality). The two keys must not be the same value.
Implementations should derive both keys from a single master secret
using a key derivation function or a deterministic hash chain that
produces cryptographically independent outputs.

\subsubsection{Collision Guard
Mechanism}\label{collision-guard-mechanism}

Without the collision guard, there exists a combined approximately
\(1/2^{48}\) probability that ciphertext coincidentally matches both the
plaintext marker bytes and produces a valid MAC when interpreted as a
plaintext SKEID. This would cause a Secure SKEID to be misclassified as
a plaintext SKEID with incorrect data.

The collision guard eliminates this edge case by construction. During
Secure SKEID generation, when SKEID Marker and MAC collision is detected
in the ciphertext, the plaintext is regenerated with successive variant
bytes from \texttt{0x8E} through \texttt{0xBF} (50 alternatives). Due to
the AES avalanche effect, each variant produces completely different
ciphertext. Termination is deterministic. The loop is bounded at 51
total attempts (1 primary + 50 alternatives), with an
implementation-defined error on exhaustion (the reference implementation
throws \texttt{JackpotException}). The probability of exhausting all 51
attempts is approximately \(1/2^{48 \times 51}\), which is vanishingly
small.

\subsubsection{Backward Verification
Algorithm}\label{backward-verification-algorithm}

During parsing, when a non-default variant byte \(V\) is recovered from
the decrypted plaintext (\(V > \text{0x8D}\)), the parse algorithm must
verify that the previous variant (\(V-1\)) genuinely triggered the
collision guard. This is achieved by reconstructing the SKEID with
variant \(V-1\), encrypting it, and checking whether the ciphertext
exhibits the SKEID Marker and MAC collision. This single-step backward
proof is sufficient by induction. If variant \(V\) is legitimate, then
\(V-1\) must have collided, and \(V-1\)'s legitimacy is either the base
case (\(V-1 = \text{0x8D}\), always legitimate) or proved by \(V-2\)
having collided (which was already verified at generation time).

The result is a deterministic guarantee. No Secure SKEID produced by a
compliant implementation can ever pass the plaintext parse path with a
valid result.

\newpage

\subsubsection{Numeric Walkthrough}\label{numeric-walkthrough}

The following example illustrates the collision guard and backward
verification with concrete values.

\textbf{Setup:} Consider a SKID with value
\texttt{0x8BEB\_C200\_1204\_0005} (timestamp = 400,000,000 ticks from
epoch, equivalent to 100,000,000 seconds at 250ms precision, App ID =
18, App Instance ID = 1, Sequence = 5), entity type \texttt{0x0A}
(entity type 10), epoch \texttt{0x00}, and key pair
\((K_{mac}, K_{aes})\).

\textbf{Step 1 SKEID Construction:} The 16-byte SKEID buffer is
constructed in big-endian order. Byte 0 receives the epoch
(\texttt{0x00}), bytes 1--4 receive the sign-toggled SKID upper half
(\texttt{0x8BEBC200} XOR \texttt{0x80000000} = \texttt{0x0BEBC200}),
byte 5 receives the most significant byte of the SKID lower half, byte 6
receives the version marker (\texttt{0x8D}), byte 7 receives the entity
type (\texttt{0x0A}), byte 8 receives the default variant marker
(\texttt{0x8D}), bytes 9--11 receive the remaining SKID lower half
bytes, and the BLAKE3 keyed MAC is computed and placed at bytes 12--15.
The UUID is constructed from the 16-byte buffer in big-endian order per
RFC 9562.

\textbf{Step 2 Encryption and Collision Check:} The 16-byte plaintext is
encrypted with AES-256-ECB using \(K_{aes}\), producing ciphertext
\(C_1\). The generator checks whether \(C_1\) coincidentally has
\texttt{0x8D} at byte position 6 and \texttt{0x8D} at byte position 8.
In the overwhelming majority of cases (probability
\(\approx 1 - 1/65{,}536\)), the ciphertext does not match, and \(C_1\)
is the final Secure SKEID.

\textbf{Step 3 Collision Scenario:} Suppose \(C_1\) happens to exhibit
\texttt{0x8D} at position 6 and \texttt{0x8D} at position 8, and
furthermore the bytes at positions 12--15 coincidentally form a valid
BLAKE3 MAC over the non-MAC bytes of \(C_1\) when interpreted as a
plaintext SKEID. This combined event has probability
\(\approx 1/2^{48}\).

When the collision guard activates, the plaintext variant byte (position
8) is changed from \texttt{0x8D} to \texttt{0x8E}. The BLAKE3 MAC is
recomputed over the modified plaintext. The new plaintext is encrypted
with the same \(K_{aes}\), producing ciphertext \(C_2\). Due to AES's
avalanche property, even a single-bit change in plaintext produces
ciphertext that differs in approximately 50\% of its bits. The
probability that \(C_2\) also triggers a collision is again
\(\approx 1/2^{48}\), independent of the first collision. In practice,
\(C_2\) passes the check and becomes the final Secure SKEID.

\textbf{Step 4 Backward Verification at Parse Time:} When a consumer
parses \(C_2\) by decrypting with \(K_{aes}\), the recovered plaintext
reveals variant byte \texttt{0x8E} (\(> \text{0x8D}\)). The parser must
verify that the escalation was legitimate.

\begin{enumerate}
\def\labelenumi{\arabic{enumi}.}
\tightlist
\item
  Reconstruct the SKEID with variant \texttt{0x8D} (replacing
  \texttt{0x8E} and recomputing the MAC).
\item
  Encrypt that reconstruction with \(K_{aes}\) to obtain \(C_1\).
\item
  Check that \(C_1\) exhibits the marker-plus-MAC coincidence,
  confirming the collision that justified the escalation.
\item
  Since \texttt{0x8D} is the base case (always legitimate), the single
  backward step completes the proof.
\end{enumerate}

If the backward check fails (the reconstructed \(C_1\) does not exhibit
the coincidence), the parser rejects the identifier as invalid. This
prevents an attacker from crafting identifiers with arbitrary variant
bytes.

Exact hex values, byte layouts, and round-trip assertions in this
walkthrough are verified by the \texttt{PaperNumericWalkthroughTests}
unit test suite in the reference implementation. Any code change that
would invalidate these values produces a test failure.

\subsection{Epoch and Extensibility}\label{epoch-and-extensibility}

The 32-bit timestamp field stores 250-millisecond ticks, giving
\(2^{32}\) ticks per epoch half (equivalent to \(2^{30}\) seconds). Each
epoch half covers approximately 34 years. The sign bit doubles this to
approximately 68 years per epoch while preserving monotonic sort order.
The sign-bit toggle in the SKEID byte layout ensures that lexicographic
comparison of UUID strings matches the signed chronological ordering of
the underlying SKIDs. SKEID byte 0 carries an 8-bit epoch index, where
each value selects a \(2^{31}\)-second window starting from 2025-01-01,
giving 256 epochs and approximately 17,421 years of total coverage.

Table 6 illustrates epoch addressing across the full epoch range.

\textbf{Table 6:} Epoch addressing.

{\def\LTcaptype{none} % do not increment counter
\begin{longtable}[]{@{}lll@{}}
\toprule\noalign{}
Epoch Value & Start & End \\
\midrule\noalign{}
\endhead
\bottomrule\noalign{}
\endlastfoot
\texttt{0x00} & January 1, 2025 & January 19, 2093 \\
\texttt{0x01} & January 19, 2093 & February 7, 2161 \\
\(\vdots\) & \(\vdots\) & \(\vdots\) \\
\texttt{0xFF} & 19378 & 19446 \\
\end{longtable}
}

This two-level time addressing gives the identity system a multi-century
lifespan without identifier collisions or sorting degradation. Epoch
boundaries, total coverage, half-epoch boundary year, and sign-bit
sort-order invariants in Table 6 are verified by the
\texttt{PaperEpochAddressabilityTests} unit test suite in the reference
implementation.

\subsection{Defense-in-Depth Security
Architecture}\label{defense-in-depth-security-architecture}

The SKEID system employs multiple verification layers. An attacker
attempting to forge a valid identifier must defeat all applicable
layers. Each layer is assumed to be independent of the others.

Table 7 enumerates the defense-in-depth security layers and their
individual bypass probabilities.

\textbf{Table 7:} Defense-in-depth security layers.

{\def\LTcaptype{none} % do not increment counter
\begin{longtable}[]{@{}
  >{\raggedright\arraybackslash}p{(\linewidth - 4\tabcolsep) * \real{0.2737}}
  >{\raggedright\arraybackslash}p{(\linewidth - 4\tabcolsep) * \real{0.4211}}
  >{\raggedright\arraybackslash}p{(\linewidth - 4\tabcolsep) * \real{0.3053}}@{}}
\toprule\noalign{}
\begin{minipage}[b]{\linewidth}\raggedright
Layer
\end{minipage} & \begin{minipage}[b]{\linewidth}\raggedright
Mechanism
\end{minipage} & \begin{minipage}[b]{\linewidth}\raggedright
Bypass Probability
\end{minipage} \\
\midrule\noalign{}
\endhead
\bottomrule\noalign{}
\endlastfoot
1. AES-256 Encryption & PRP over 128 bits & \(2^{-128}\) without the
key \\
2. BLAKE3 Keyed MAC & 32-bit truncated & \(2^{-32}\) per attempt \\
3. Marker Detection & \texttt{0x8D8D} at fixed positions &
\(2^{-16}\) \\
4. Entity Type Match & Must match the expected type & \(2^{-8}\) \\
5. Topology Validity & App ID and Instance ID must exist & Variable \\
6. Existence Probability & Forged SKID must reference a record &
Variable \\
7. Rate Limiting & Restricts attempts per time period & Implementation
defined \\
\end{longtable}
}

Even if an attacker guesses a ciphertext that, when decrypted, has valid
markers, a valid MAC, a valid entity type, and valid topology, the
resulting SKID must still correspond to an actual record in the
database. The multiplicative combination of these barriers offers
security far exceeding any individual layer. The layers are assumed to
be operationally independent under secure key derivation. Each layer
uses a different cryptographic primitive or validation mechanism. A
formal proof of probabilistic independence is outside the scope of this
paper (see Limitation 4).

\subsubsection{MAC Truncation Analysis}\label{mac-truncation-analysis}

The 32-bit truncation satisfies the minimum MacTag length specified in
NIST SP 800-107 (Dang, 2012). The document notes that MacTags shorter
than 64 bits are discouraged in Section 5.3.5. The 128-bit SKEID format
constrains the available space for the MAC field, as the remaining bits
carry the SKID, entity type, epoch, and marker bytes. This constraint is
acceptable because the MAC is not the sole security mechanism. It
operates within the defense-in-depth architecture described in Table 7,
where AES-256 encryption, marker detection, entity type matching, and
record existence probability provide additional verification layers.

Section 5.3.5 of the same document provides the general framework for
assessing truncated MAC forgery risk. For a \(\lambda\)-bit MacTag with
\(2^t\) failed verifications allowed, the likelihood of accepting forged
data is \((1/2)^{(\lambda - t)}\). Applying this to the 32-bit SKEID MAC
(\(\lambda = 32\)): if a system permits \(2^{12}\) (4,096) failed
verification attempts before rotating the MAC key, the forgery
likelihood is \((1/2)^{20}\), approximately one in a million. At 100
rate-limited attempts per second, this \(2^{12}\) budget is exhausted in
roughly 41 seconds.

SP 800-107 Rev.~1 was withdrawn in 2022 (NIST, 2022). Its successor,
NIST SP 800-224 (Turan \& Brandão, 2024), preserves the same truncated
MAC analysis. SP 800-224 stipulates that tag lengths below 64 bits
require careful risk analysis, which this section and the
defense-in-depth architecture in Table 7 do.

\subsection{Key Rotation and Key-Ring
Fallback}\label{key-rotation-and-key-ring-fallback}

A key-ring containing one or more key entries is required. Exactly one
key is designated as the default (active) key at any time. Previous keys
remain in the key-ring for parse fallback. All new SKEID and Secure
SKEID generation uses the current default key. When parsing, the system
attempts verification with the current default key first. If
verification fails, it iterates through previous keys in reverse
chronological order.

Key rotation does not require re-encrypting existing database records
because the 64-bit SKID stored in the database contains no cryptographic
material. Given the SKID and entity type, a new SKEID or Secure SKEID
can be regenerated with the new key at any time. In the event of key
compromise:

\begin{enumerate}
\def\labelenumi{\arabic{enumi}.}
\tightlist
\item
  Add the compromised key to the key-ring with a ``compromised'' label.
\item
  Set a new key as the default.
\item
  Optionally regenerate SKEIDs from the immutable SKIDs.
\end{enumerate}

\subsection{AI-Assisted Research and
Development}\label{ai-assisted-research-and-development}

\begin{quote}
This subsection fulfills PeerJ CS's AI disclosure requirement for AI
tools used as part of the research methodology.
\end{quote}

The following AI tools and large language model (LLM) systems were used
in the development of the reference implementation and the preparation
of this research.

\textbf{Software development, documentation, and testing:} AI tools from
Anthropic (Claude Opus/Sonnet series), Google (Gemini Pro/Flash series),
OpenAI (o1, 4o), Qwen 3, and DeepSeek R1 were used across the
DRN-Project development lifecycle, including the SKIDs implementation,
as agentic coding assistants for code generation, refactoring, code
review, test scaffolding, and documentation drafting. SKIDs were added
to the DRN-Project roadmap on 19 November 2023. These AI tools were
progressively adopted as they became available over the subsequent
two-plus years of development. All generated code, tests, and
documentation were reviewed, validated, and modified by the author.

\textbf{Literature research:} Google Gemini Deep Research was used to
assist with surveying existing distributed identifier schemes and
locating relevant prior work. All identified sources were independently
verified by the author before citation.

\textbf{Manuscript preparation:} Claude Opus 4.6 and Gemini 3.1 Pro were
used to draft, structure, and refine sections of this paper. All
technical claims were verified against the implementation, and final
editorial decisions were made by the author.

\textbf{Prompt disclosure:} AI interactions followed the following
agentic workflow pattern. The author provided high-level task
descriptions, architectural constraints, and acceptance criteria. The AI
tools generated candidate implementations, which the author reviewed,
modified, and validated. Following representative prompt categories
included.

\begin{itemize}
\tightlist
\item
  ``implement Secure SKEID generation with the following byte layout by
  using AES-256-ECB''
\item
  ``refactor sequence manager to use lock-free atomic operations''
\item
  ``write integration tests for Secure SKEID round-trip encryption using
  Testcontainers''
\item
  ``review this section for technical accuracy against the reference
  implementation and explain your judgement''
\end{itemize}

Due to the iterative, multi-session nature of the two-year development
effort across multiple AI tools, exhaustive prompt logs are not
available in a reproducible form. The prompts described above represent
the general interaction patterns used throughout the project.

AI-assisted development workflows were guided by a structured behavioral
framework enforcing security-first principles and systematic
decision-making.

All AI-assisted outputs were reviewed, validated, and refined by the
author. The architectural design, cryptographic choices, three-tier
trust model, and all technical decisions are solely the work of the
author. The author takes full responsibility for the correctness and
originality of the work.

\newpage

\subsection{Reference Implementation}\label{reference-implementation}

The reference implementation is provided as part of the open-source
DRN.Framework (Kılıç, 2023), organized across four NuGet packages.

\begin{itemize}
\tightlist
\item
  \textbf{DRN.Framework.SharedKernel} defines interfaces
  (\texttt{ISourceKnownEntityIdOperations}), base classes
  (\texttt{SourceKnownEntity}), and parsed identity structs
  (\texttt{SourceKnownId}, \texttt{SourceKnownEntityId}).
\item
  \textbf{DRN.Framework.Utils} implements cryptographic operations,
  timestamp management, sequence management
  (\texttt{SequenceManager\textless{}TEntity\textgreater{}} with
  lock-free atomic operations), and bit packing
  (\texttt{SourceKnownIdUtils}, \texttt{SourceKnownEntityIdUtils}).
\item
  \textbf{DRN.Framework.EntityFramework} implements persistence
  integration including \texttt{SourceKnownIdValueGenerator} (EF Core
  value generator for automatic SKID assignment),
  \texttt{DrnSaveChangesInterceptor} (auto-assigns SKID and SKEID on
  insert), \texttt{DrnMaterializationInterceptor} (reconstructs full
  SKEID from stored SKID on database read), and
  \texttt{SourceKnownRepository\textless{}TContext,\ TEntity\textgreater{}}
  (generic repository with SKEID validation and cursor-based
  pagination).
\item
  \textbf{DRN.Framework.Testing} implements integration and unit test
  infrastructure that provides Testcontainers orchestration for
  ephemeral PostgreSQL instances, auto-mocking through data attributes
  (\texttt{{[}DataInline{]}}, \texttt{{[}DataInlineUnit{]}}), and
  convention-based test contexts (\texttt{DrnTestContext},
  \texttt{DrnTestContextUnit}) for low-friction, high-fidelity testing.
\end{itemize}

The implementation uses attribute-based dependency injection
(\texttt{{[}Singleton\textless{}Type\textgreater{}{]}},
\texttt{{[}Scoped\textless{}Type\textgreater{}{]}},
\texttt{{[}Transient\textless{}Type\textgreater{}{]}}) and entity type
attributes (\texttt{{[}EntityType(byte){]}} for entity type
registration) and integrates with Domain-Driven Design (Evans, 2003)
patterns through the \texttt{SourceKnownEntity} abstract base class. The
testing strategy and validation evidence for these packages are
described in the Validation and Maturity section.

\section{Results}\label{results}

\subsection{Reference Benchmark
Environment}\label{reference-benchmark-environment}

All benchmarks were conducted using BenchmarkDotNet ({Akinshin et al.},
2013) v0.15.8 on the following hardware and software configuration.

\begin{itemize}
\tightlist
\item
  \textbf{Hardware:} Apple M2, 1 CPU, 8 logical and 8 physical cores
\item
  \textbf{Operating System:} macOS Tahoe 26.4 (25E246) {[}Darwin
  25.4.0{]}
\item
  \textbf{Runtime:} .NET 10.0.5 (10.0.526.15411), Arm64 RyuJIT armv8.0-a
\item
  \textbf{SDK:} .NET SDK 10.0.201
\item
  \textbf{Configuration:} OutlierMode=RemoveUpper,
  InvocationCount=262,144, IterationCount=120, WarmupCount=120
\end{itemize}

OutlierMode is set to RemoveUpper, a standard BenchmarkDotNet
configuration that removes upper statistical outliers caused by garbage
collector (GC) pauses and operating system (OS) scheduling jitter before
computing summary statistics. This isolates measurements to the
intrinsic operation cost.

The benchmark configuration uses 120 iterations with 262,144 invocations
per iteration, enabling BenchmarkDotNet to compute error margins and
standard deviations that approximate normal distribution. The invocation
count of 262,144 matches the SKID system's 18-bit sequence counter cap
(\(2^{18}\) identifiers per 250ms tick per instance). All invocations
within a single iteration complete within the 250-millisecond sequence
time scope, isolating the intrinsic computational cost of each operation
without triggering the backpressure mechanism. A 250-millisecond
\texttt{IterationSetup} delay between iterations allows the sequence
time scope to reset, so each iteration starts from a clean state.
Separate saturation benchmarks with \texttt{InvocationCount=786,432} (3
times the sequence cap), \texttt{IterationCount=40}, and
\texttt{WarmupCount=40} (matching the iteration count) confirm the
backpressure behavior. The saturation benchmark uses BenchmarkDotNet's
default OutlierMode, retaining all measurements including upper outliers
to capture the full backpressure distribution. Both configurations
produce the same total invocation count
(\(120 \times 262{,}144 = 40 \times 786{,}432 = 31{,}457{,}280\)). The
saturation benchmark is configured to observe backpressure behavior
rather than to produce directly comparable absolute timings. SKID
generation under saturation shows elevated mean times, reflecting the
intermittent wait-for-next-tick pauses rather than the true
per-operation cost. This validates the design. The sequence counter
applies backpressure to preserve uniqueness guarantees under sustained
load.

\subsection{Performance Results}\label{performance-results}

Table 8 presents the BenchmarkDotNet performance measurements for all
SKID system operations. Mean, Error, and Standard Deviation values are
reported in nanoseconds (ns) and rounded to a maximum of two decimal
places. Error values represent the 99.9\% CI margin computed by
BenchmarkDotNet over iterations remaining after upper outlier removal.

\textbf{Table 8:} BenchmarkDotNet performance measurements (120
configured iterations, 262,144 invocations per iteration).

{\def\LTcaptype{none} % do not increment counter
\begin{longtable}[]{@{}
  >{\raggedright\arraybackslash}p{(\linewidth - 8\tabcolsep) * \real{0.4500}}
  >{\raggedright\arraybackslash}p{(\linewidth - 8\tabcolsep) * \real{0.1200}}
  >{\raggedright\arraybackslash}p{(\linewidth - 8\tabcolsep) * \real{0.1200}}
  >{\raggedright\arraybackslash}p{(\linewidth - 8\tabcolsep) * \real{0.1300}}
  >{\raggedright\arraybackslash}p{(\linewidth - 8\tabcolsep) * \real{0.1800}}@{}}
\toprule\noalign{}
\begin{minipage}[b]{\linewidth}\raggedright
Operation
\end{minipage} & \begin{minipage}[b]{\linewidth}\raggedright
Mean (ns)
\end{minipage} & \begin{minipage}[b]{\linewidth}\raggedright
Error (ns)
\end{minipage} & \begin{minipage}[b]{\linewidth}\raggedright
StdDev (ns)
\end{minipage} & \begin{minipage}[b]{\linewidth}\raggedright
Allocated Memory
\end{minipage} \\
\midrule\noalign{}
\endhead
\bottomrule\noalign{}
\endlastfoot
Random 64-bit integer generation & 171.3 & 2.36 & 7.66 & 32 B \\
Random UUID V4 generation & 354.1 & 3.64 & 11.81 & 0 B \\
Random UUID V7 generation & 377.5 & 3.24 & 10.53 & 0 B \\
Timestamp manager current time generation & 4.1 & 0.46 & 1.43 & 0 B \\
Sequence manager time-scoped ID generation & 15.4 & 0.69 & 2.19 & 0 B \\
\textbf{SKID generation} & \textbf{35.3} & \textbf{1.21} & \textbf{3.92}
& \textbf{0 B} \\
SKEID generation (with provided SKID) & 219.0 & 3.25 & 10.56 & 0 B \\
\textbf{SKEID generation (with SKID generation)} & \textbf{230.3} &
\textbf{3.45} & \textbf{11.20} & \textbf{0 B} \\
SKEID generation (with entity allocation) & 248.3 & 3.87 & 12.58 & 192
B \\
\textbf{Secure SKEID generation} & \textbf{544.0} & \textbf{5.67} &
\textbf{18.42} & \textbf{72 B} \\
SKEID parsing & 223.6 & 2.87 & 9.31 & 0 B \\
Secure SKEID parsing & 540.7 & 3.15 & 10.22 & 72 B \\
ToPlain (decryption only) & 217.3 & 2.40 & 7.76 & 0 B \\
ToSecure (encryption only) & 524.2 & 5.56 & 17.98 & 72 B \\
\end{longtable}
}

\subsection{Performance Comparison with
UUID}\label{performance-comparison-with-uuid}

Secure SKEID generation at 544.0 ns takes approximately 1.4 times as
long as UUID V7 at 377.5 ns. This overhead adds AES-256 encryption,
BLAKE3 MAC integrity, entity type discrimination, and zero-lookup
verification, none of which are provided by UUID V7. Secure SKEID can be
converted to SKEID or SKID for the other desired identifier properties.
The additional cost reflects approximately 314 ns of AES-256
single-block encryption beyond plaintext SKEID generation.

SKEID generation at 230.3 ns is approximately 1.6 times as fast as UUID
V7 (377.5 ns) despite embedding additional metadata (entity type, epoch)
and computing a BLAKE3 keyed MAC within the same 128-bit footprint. In
the reference implementation, SKEID generation operates entirely on the
stack (0 bytes allocated), avoiding garbage collection pressure.
Implementations in other managed runtimes may exhibit different
allocation profiles depending on stack allocation support.

SKID generation at 35.3 ns is more than 10 times as fast as UUID V7
generation at 377.5 ns and more than 10 times as fast as UUID V4 at
354.1 ns. This performance advantage stems from the deterministic
bit-packing approach. SKID generation requires only an atomic counter
increment, a cached timestamp read, and bitwise operations, with no
cryptographic random number generation.

\subsection{Throughput Analysis}\label{throughput-analysis}

The 18-bit sequence field caps generation at 262,144 identifiers per
250ms tick (1,048,576 per second) per instance. With 128 applications
and 64 instances per application (8,192 total generators), the
theoretical maximum system-wide throughput is approximately 8.6 billion
identifiers per second. Full-throttle benchmarks confirm that the
sequence manager applies backpressure when the per-instance limit is
reached, preserving uniqueness under sustained load. Under saturation
conditions (786,432 invocations per iteration), SKID generation mean
rises from 35.3 ns to 610.0 ns because two-thirds of invocations exceed
the per-tick sequence cap and must wait for the next 250-millisecond
tick boundary. The saturation benchmark confirms that backpressure is
precisely scoped to the sequence manager. All operations that internally
generate a SKID (SKEID generation, Secure SKEID generation) show
elevated means consistent with tick-boundary waiting, while operations
that do not generate a SKID (parsing, tier conversion with a provided
identifier) remain unaffected. This isolation validates the
architectural separation between identifier generation and identifier
processing (see supplementary file
\texttt{supplementary-saturation-report.html}).

Throughput limits, generator topology, and system-wide capacity in this
section are verified by the \texttt{PaperThroughputAnalysisTests} unit
test suite in the reference implementation.

\subsection{Storage Comparison}\label{storage-comparison}

Table 9 compares storage requirements across identifier schemes.

\textbf{Table 9:} Storage requirements comparison.

{\def\LTcaptype{none} % do not increment counter
\begin{longtable}[]{@{}
  >{\raggedright\arraybackslash}p{(\linewidth - 6\tabcolsep) * \real{0.2895}}
  >{\raggedright\arraybackslash}p{(\linewidth - 6\tabcolsep) * \real{0.2632}}
  >{\raggedright\arraybackslash}p{(\linewidth - 6\tabcolsep) * \real{0.2632}}
  >{\raggedright\arraybackslash}p{(\linewidth - 6\tabcolsep) * \real{0.1842}}@{}}
\toprule\noalign{}
\begin{minipage}[b]{\linewidth}\raggedright
Scheme
\end{minipage} & \begin{minipage}[b]{\linewidth}\raggedright
Primary Key Size
\end{minipage} & \begin{minipage}[b]{\linewidth}\raggedright
External ID Size
\end{minipage} & \begin{minipage}[b]{\linewidth}\raggedright
Ordered
\end{minipage} \\
\midrule\noalign{}
\endhead
\bottomrule\noalign{}
\endlastfoot
SKID/SKEID & 8 B & 16 B & Yes \\
UUID V4 & 16 B & 16 B & No \\
UUID V7 & 16 B & 16 B & Yes \\
Snowflake & 8 B & 8 B & Yes \\
ULID & 16 B & 26 chars (Base32) & Yes \\
CUID2 & 24 chars & 24 chars & No \\
KSUID & 20 B & 27 chars (Base62) & Yes \\
DB Sequence & 4/8 B & 4/8 B & Yes \\
Dual ID (int + UUID) & 8 B + 16 B = 24 B & 16 B & Integer only \\
\end{longtable}
}

At 8 bytes per SKID versus 16 bytes for UUID, storage cost halves,
directly reducing archival and indexing overhead. Compared to KSUIDs at
20 bytes, storage cost is reduced by a factor of 2.5 per primary key.
SKIDs as 64-bit integers provide superior B-tree index performance
because smaller keys mean more keys per B-tree page, fewer page splits,
and better cache utilization. The timestamp-leading layout means new
records append to the end of the index, optimizing for append-heavy
workloads. Cursor-based pagination is natively supported through the
SKID's monotonic nature
(\texttt{WHERE\ id\ \textgreater{}\ :lastId\ ORDER\ BY\ id\ LIMIT\ :pageSize}),
eliminating offset-based pagination overhead. The dual-identifier
pattern (integer primary key plus UUID external identifier) requires
maintaining two separate columns with 24 bytes total per record. SKIDs
eliminate this redundancy. The 8-byte SKID serves as the primary key,
and the 16-byte SKEID or Secure SKEID is computed deterministically on
demand without additional storage.

The computational cost of tier transformations is negligible. Converting
between representations requires only sub-microsecond CPU operations.
ToSecure (encryption) completes in 524.2 ns and ToPlain (decryption) in
217.3 ns (Table 8), with no I/O, no network round-trip, and no external
dependency. Any alternative that relies on a secondary lookup, whether a
database join for the dual-identifier pattern, a cache query, or a
remote service call, incurs latency orders of magnitude higher. This
makes the transformation cost effectively zero relative to any I/O-bound
identifier resolution strategy.

\newpage

\section{Discussion}\label{discussion}

\subsection{Interpretation of Results}\label{interpretation-of-results}

Benchmark results of the reference implementation demonstrate that the
SKID system achieves competitive or superior performance compared to
standard identifier generation while delivering substantially more
functionality. The more than 10× speed advantage of SKID over UUID V7
stems from the deterministic generation approach. Bit packing from
cached values is fundamentally less expensive than the cryptographic
random number generation required by UUID schemes. Even the most
expensive operation (Secure SKEID at 544.0 ns) completes in well under a
microsecond, making it practical for high-throughput production systems.

Zero-allocation behavior in core SKID and SKEID operations is important
for managed runtime environments (e.g., .NET, JVM) where garbage
collection pauses can impact tail latency. By operating entirely on the
stack, SKID generation avoids contributing to GC pressure even under
sustained high-throughput scenarios.

Throughput scales linearly with the number of active generators. A
single generator instance produces up to 1,048,576 identifiers per
second (262,144 per 250ms tick times 4 ticks per second), bounded by the
18-bit sequence counter. With the default topology configuration
supporting 8,192 concurrent generators (128 applications times 64
instances), the theoretical system-wide ceiling is approximately 8.6
billion identifiers per second. Most distributed applications will not
approach this ceiling. The backpressure mechanism ensures that
overloaded generators degrade gracefully by waiting for the next tick
rather than producing duplicates, a deterministic safety guarantee that
probabilistic schemes cannot offer.

\subsection{Comparison with Existing
Work}\label{comparison-with-existing-work}

Compared to UUID V7 (Davis, Peabody \& Leach, 2024), the most directly
comparable recent standard, SKIDs offer origin metadata, integrity
verification (MAC), and confidentiality (encryption) that UUID V7 lacks.
UUID V7 provides millisecond-precision timestamps while SKID uses
250-millisecond tick precision (see Design Trade-offs for the rationale
and throughput analysis of this choice).

Snowflake (Twitter Engineering, 2010) shares the same 64-bit compact
footprint as SKID, but the SKEID extension layer adds integrity and
confidentiality that Snowflake lacks. Snowflake's 10-bit machine ID
(1,024 machines) is comparable to SKID's combined 13-bit topology (7-bit
App ID + 6-bit Instance ID), but SKEID adds explicit entity type
discrimination. Both require coordination to assign machine/instance
identifiers and both schemes converge on similar single-epoch durations
(\textasciitilde68 years for SKID, \textasciitilde69.7 years for
Snowflake). However, Snowflake's epoch is non-extensible while SKEID can
extend SKID with an 8-bit epoch index that provides 256 configurable
epochs for approximately 17,421 years of total coverage. Lastly,
Snowflake uses 41-bit millisecond timestamps while SKID achieves
comparable coverage with 32-bit 250-millisecond ticks.

ULIDs (Feerasta, 2016) require twice the primary key storage (16 bytes
versus 8 for SKID) and carry less embedded metadata. ULIDs use a 48-bit
millisecond timestamp with \textasciitilde8,919 years of coverage from
the Unix epoch, while SKIDs achieve \textasciitilde68 years per epoch
with 256 configurable epochs for \textasciitilde17,421 years total.
ULIDs rely on 80 bits of cryptographic randomness for uniqueness within
the same millisecond, a probabilistic guarantee, while SKIDs use
deterministic topology-based partitioning. ULIDs use a 26-character
Crockford Base32 encoding that sorts lexicographically, targeting
string-first environments, but lack UUID format compatibility. SKEIDs
are valid RFC 9562 Version 8 UUIDs. ULIDs provide neither integrity
verification nor confidentiality.

KSUIDs (Segment, 2017) are the largest scheme in this comparison at 20
bytes per identifier (see Table 9). SKIDs provide a more compact
representation (8/16 bytes) with richer metadata. KSUIDs use
second-precision timestamps while SKIDs use 250-millisecond ticks.
KSUIDs provide \textasciitilde136 years coverage with 32-bit seconds and
no extension mechanism. SKIDs provide \textasciitilde68 years per epoch
with 256 configurable epochs for \textasciitilde17,421 years total.
KSUID's 128-bit random payload trades storage efficiency for stronger
per-identifier collision resistance without coordination. SKIDs achieve
uniqueness through the topology-based partitioning scheme with
deployment time coordination, which guarantees uniqueness by
construction rather than by statistical improbability. Additionally,
SKIDs provide entity type discrimination, integrity verification, and
confidentiality, none of which KSUID offers. Even at the external-facing
tier, SKEIDs (16 bytes) remain more compact than KSUIDs (20 bytes) while
satisfying all six desired identifier properties. KSUIDs also lack UUID
format compatibility, using a custom 27-character Base62 encoding that
requires specialized tooling rather than standard UUID libraries
available in every major programming language.

The original CUID (Elliott, 2012) was deprecated due to the security
vulnerabilities inherent in k-sortable, timestamp-prefixed identifiers,
as discussed in the Related Work section. Secure SKEIDs address this
class of vulnerability directly. The AES-256 encryption layer transforms
the entire UUID V8 payload into ciphertext indistinguishable from random
bytes, preventing the information leakage that the CUID deprecation
notice identifies as a systemic risk across timestamp-based and
k-sortable identifier schemes.

CUID2 (Elliott, 2022) achieves uniqueness through SHA3-hashed entropy, a
probabilistic guarantee. SKIDs instead offer deterministic uniqueness by
construction. The CUID2 authors observe that ``insecure ids can cause
problems in unexpected ways, including unauthorized user account access,
unauthorized access to user data, and accidental leaks of user's
personal data which can lead to catastrophic effects, even in
innocent-sounding applications like fitness run trackers'' (Elliott,
2022). CUID2 addressed the security issues of its predecessor by
replacing the k-sortable structure with SHA3-hashed entropy. CUID2
generates identifiers by combining multiple entropy sources and hashing
them with SHA3, producing strong collision resistance but relying on the
statistical properties of the hash function rather than structural
guarantees. CUID2 also recommends a separate \texttt{createdAt} database
column for time-based sorting. That change introduces additional
per-record storage overhead and a second indexed column, a trade-off
structurally similar to the dual-identifier pattern (integer primary key
plus UUID) that SKIDs eliminate by embedding the timestamp directly in
the identifier.

SKIDs prefer context-dependent uniqueness over probabilistic uniqueness.
This approach provides efficiency without sacrificing uniqueness
guarantees.

Tables 10 and 11 compare the security posture of the SKEID system
against standard and alternative identifier schemes.

\textbf{Table 10:} Security comparison of standard identifier schemes.

{\def\LTcaptype{none} % do not increment counter
\begin{longtable}[]{@{}
  >{\raggedright\arraybackslash}p{(\linewidth - 8\tabcolsep) * \real{0.2178}}
  >{\raggedright\arraybackslash}p{(\linewidth - 8\tabcolsep) * \real{0.2178}}
  >{\raggedright\arraybackslash}p{(\linewidth - 8\tabcolsep) * \real{0.1683}}
  >{\raggedright\arraybackslash}p{(\linewidth - 8\tabcolsep) * \real{0.1881}}
  >{\raggedright\arraybackslash}p{(\linewidth - 8\tabcolsep) * \real{0.2079}}@{}}
\toprule\noalign{}
\begin{minipage}[b]{\linewidth}\raggedright
Threat
\end{minipage} & \begin{minipage}[b]{\linewidth}\raggedright
SKEID / Secure SKEID
\end{minipage} & \begin{minipage}[b]{\linewidth}\raggedright
UUID V4
\end{minipage} & \begin{minipage}[b]{\linewidth}\raggedright
UUID V7
\end{minipage} & \begin{minipage}[b]{\linewidth}\raggedright
Snowflake
\end{minipage} \\
\midrule\noalign{}
\endhead
\bottomrule\noalign{}
\endlastfoot
Enumeration & MAC + optional AES & 122-bit random & Timestamp-ordered &
Timestamp-ordered \\
Forgery & MAC verification & No detection & No detection & No
detection \\
Info leakage & AES-256 encrypted & None (random) & Timestamp visible &
TS + worker visible \\
Cross-type confusion & Entity type check & Possible & Possible &
Possible \\
\end{longtable}
}

\newpage

\textbf{Table 11:} Security comparison of alternative identifier
schemes.

{\def\LTcaptype{none} % do not increment counter
\begin{longtable}[]{@{}
  >{\raggedright\arraybackslash}p{(\linewidth - 8\tabcolsep) * \real{0.2178}}
  >{\raggedright\arraybackslash}p{(\linewidth - 8\tabcolsep) * \real{0.2178}}
  >{\raggedright\arraybackslash}p{(\linewidth - 8\tabcolsep) * \real{0.1980}}
  >{\raggedright\arraybackslash}p{(\linewidth - 8\tabcolsep) * \real{0.1584}}
  >{\raggedright\arraybackslash}p{(\linewidth - 8\tabcolsep) * \real{0.2079}}@{}}
\toprule\noalign{}
\begin{minipage}[b]{\linewidth}\raggedright
Threat
\end{minipage} & \begin{minipage}[b]{\linewidth}\raggedright
SKEID / Secure SKEID
\end{minipage} & \begin{minipage}[b]{\linewidth}\raggedright
ULID
\end{minipage} & \begin{minipage}[b]{\linewidth}\raggedright
CUID2
\end{minipage} & \begin{minipage}[b]{\linewidth}\raggedright
KSUID
\end{minipage} \\
\midrule\noalign{}
\endhead
\bottomrule\noalign{}
\endlastfoot
Enumeration & MAC + optional AES & TS + 80-bit random & Hash-based & TS
+ 128-bit random \\
Forgery & MAC verification & No detection & No detection & No
detection \\
Info leakage & AES-256 encrypted & Timestamp visible & None (hashed) &
Timestamp visible \\
Cross-type confusion & Entity type check & Possible & Possible &
Possible \\
\end{longtable}
}

\subsection{Threat Analysis}\label{threat-analysis}

To evaluate the security posture of the SKEID system systematically, we
apply a threat analysis informed by the STRIDE categories (Shostack,
2014) (Spoofing, Tampering, Repudiation, Information Disclosure, Denial
of Service, Elevation of Privilege) to the three-tier identity model.

\textbf{Spoofing (Identifier Fabrication):} An attacker who intercepts a
Secure SKEID cannot derive the plaintext SKEID without the AES-256 key.
Forging a valid Secure SKEID requires producing ciphertext that, when
decrypted, yields valid markers, a correct BLAKE3 MAC, a valid entity
type, and a SKID corresponding to an existing record. The multiplicative
barrier across layers 1--6 in Table 7 makes this computationally
infeasible. For plaintext SKEIDs within trusted environments, the BLAKE3
MAC prevents identifier fabrication without the MAC key.

\textbf{Tampering (Identifier Modification):} Any modification to a
Secure SKEID produces different plaintext upon decryption (AES-256 PRP
property), invalidating the MAC with probability \(1 - 2^{-32}\). For
plaintext SKEIDs, modifying any byte invalidates the MAC because the MAC
is computed over the full 16-byte buffer.

\textbf{Repudiation:} The SKEID system does not provide non-repudiation
in the cryptographic sense because it does not use digital signatures.
However, the embedded topology metadata (App ID, App Instance ID, Entity
Type) and timestamp provide forensic traceability. For a given SKEID,
the originating application, instance, entity type, and approximate
creation time can be determined.

\textbf{Information Disclosure:} Plaintext SKEIDs deliberately expose
metadata (timestamp, entity type, topology) within trusted environments
where this information supports routing and validation. For external
consumers, Secure SKEIDs encrypt the entire identifier using
AES-256-ECB, which functions as a pseudorandom permutation on the single
128-bit block. The ciphertext is computationally indistinguishable from
random bytes without the key, preventing timestamp analysis, entity type
inference, and generation pattern discovery.

\textbf{Denial of Service:} The 18-bit sequence counter limits
generation to 262,144 identifiers per 250ms tick (1,048,576 per second)
per instance. Exceeding this rate triggers backpressure (the generator
waits for the next tick), which is a deliberate safety mechanism rather
than a vulnerability. An attacker cannot cause identifier exhaustion
across the system because each of the 8,192 possible generators operates
independently. Parse operations (both encrypted and plaintext) are
bounded-time with no external dependencies, preventing
parse-amplification attacks.

\textbf{Elevation of Privilege (Cross-Entity Confusion):} The 8-bit
entity type discriminator prevents cross-entity attacks where an
identifier for one entity type (e.g., User with entity type 1) is
submitted as another (e.g., AdminRole with entity type 5), potentially
granting unintended access or administrative privileges. Because the MAC
includes the entity type in its computation, a valid SKEID for entity
type 1 cannot pass MAC verification when checked against entity type 5.

\newpage

\subsection{Validation and Maturity}\label{validation-and-maturity}

Active development of the SKID system began in November 2023 as part of
the DRN.Framework. The framework has undergone continuous refinement
across versioned releases (v0.1.0 through v0.8.1). Each release is
validated through integration and unit tests built on
DRN.Framework.Testing. Performance is measured separately with
BenchmarkDotNet.

The framework's correctness is validated through multiple test and
analysis tiers.

\begin{itemize}
\tightlist
\item
  \textbf{Unit tests:} In-memory tests covering SKID generation, SKEID
  construction with MAC verification, and Secure SKEID
  encryption/decryption round-trips.
\item
  \textbf{Paper-verification tests:} Dedicated suites
  (\texttt{PaperNumericWalkthroughTests},
  \texttt{PaperEpochAddressabilityTests}) assert the exact numeric
  values, byte layouts, epoch boundaries, lexicographic ordering across
  epoch halves, and sort-order invariants presented in this paper. Any
  implementation change that would invalidate a paper claim produces a
  test failure.
\item
  \textbf{Integration tests:} Testcontainers-based tests provision
  ephemeral PostgreSQL instances, exercising Entity Framework Core value
  generation, interceptors, and cursor-based pagination with SKID
  ordering.
\item
  \textbf{End-to-end validation:} Test applications validated behavior
  across the three tiers: SKID assignment during entity creation, SKEID
  exposure within trusted application-to-application communication, and
  Secure SKEID generation for external API responses. These applications
  confirm that bidirectional transformations preserve data integrity
  across the complete round-trip.
\item
  \textbf{Static analysis:} GitHub CodeQL and SonarCloud continuously
  scan the codebase for security vulnerabilities, code quality issues,
  and potential bugs. The framework maintains zero critical or
  high-severity findings.
\end{itemize}

The four NuGet packages have been published on NuGet.org since the
initial framework release and are available for community review and
adoption. They are published through an automated GitHub Actions CI/CD
pipeline that runs the complete test suite before each release.
Therefore, no regression reaches published consumers.

\subsection{Design Trade-offs}\label{design-trade-offs}

\textbf{250-millisecond tick versus millisecond precision:} SKIDs use
250-millisecond tick timestamps (32 bits, four ticks per second) instead
of millisecond-precision (48 bits in UUID V7). This sub-second
quantization provides finer ordering than second-precision schemes while
preserving sufficient bits for application topology and sequence fields
within 64 bits. The 18-bit sequence (262,144 identifiers per tick;
1,048,576 per second) delivers high throughput within each tick.
Customized domain-tailored bit layouts can be more suitable for specific
use cases.

\textbf{Clock synchronization versus drift protection:} The SKID system
delegates clock synchronization to the operating system. Accurate
timekeeping through Network Time Protocol (NTP) (Mills et al., 2010),
Precision Time Protocol (PTP) (IEEE, 2020), or equivalent protocols is
an infrastructure responsibility outside the scope of the identifier
generator. The clock drift protection mechanism described in the SKID
specification handles only the consequences of clock jumps. Minor
backward drifts freeze the timestamp until the wall clock catches up,
while drifts beyond the freeze threshold force an application restart
with a new instance identifier. This separation of concerns keeps the
identifier generator simple and deterministic, relying on the operating
system to maintain monotonic time accuracy.

\textbf{32-bit truncated MAC versus full MAC:} The 4-byte MAC is shorter
than typical MACs (16--32 bytes) but is one layer in a defense-in-depth
strategy. The MAC is not the sole security mechanism. It works in
concert with AES encryption, marker detection, entity type matching, and
record existence probability.

\textbf{Coordination versus probabilistic uniqueness:} The SKID system
requires deployment-time coordination to assign \texttt{AppId} (7 bits,
up to 128 applications) and \texttt{AppInstanceId} (6 bits, up to 64
instances per application). This stands in contrast to UUID V4, UUID V7,
and CUID2, which achieve uniqueness without any coordination through
cryptographic random number generation. The coordination cost is an
intentional architectural trade-off. Deterministic uniqueness by
construction eliminates the residual collision probability inherent in
probabilistic schemes, enables zero-lookup verification through MAC and
topology metadata, and allows the compact 64-bit representation that
probabilistic schemes cannot achieve at equivalent uniqueness
guarantees.

\textbf{Topology field sizing:} The 7-bit application field (128
applications) and 6-bit instance field (64 instances per application)
reflect practical operational boundaries rather than arbitrary bit
allocation. Beyond 128 independently deployed applications, coordination
complexity becomes infeasible due to the increase in the number of
intercommunication paths (\(n(n-1)/2\)) (Jr., 1995, chap. 2).
Well-designed systems mitigate this through bounded contexts (Evans,
2003) rather than unbounded application proliferation. Similarly, beyond
approximately 10 concurrently active instances per application, resource
contention on shared infrastructure (databases, message brokers)
typically becomes the throughput bottleneck before the identifier
generator itself. The 6-bit field accommodates up to 64 instance
identifiers, providing headroom for operational needs such as rolling
deployments and application restarts triggered by clock drift
protection, where the restarting instance must acquire a new instance
identifier to prevent duplicate identifiers. These limits are defaults
for the general case. The system is intentionally optimized for stable,
bounded distributed application topologies. Domain-tailored bit layouts
(see Future Work, item 3) remain available for specialized deployments
requiring different trade-off profiles.

\textbf{.NET Guid mixed endianness:} The .NET \texttt{Guid} struct uses
a mixed-endian internal representation. The first three fields
(\texttt{Data1}, \texttt{Data2}, \texttt{Data3}) are stored in
little-endian order, while the remaining eight bytes (\texttt{Data4})
are stored sequentially. This internal layout is an implementation
artifact of the Windows COM heritage and differs from the big-endian
(network byte order) specified by RFC 9562. The SKEID implementation
uses the .NET 8+ constructors
\texttt{Guid(ReadOnlySpan\textless{}byte\textgreater{},\ bigEndian:\ true)}
and \texttt{Guid.ToByteArray(bigEndian:\ true)} to ensure that byte
arrays are always constructed and interpreted in RFC 9562 big-endian
format, regardless of the platform's internal \texttt{Guid}
representation. The string representation produced by
\texttt{Guid.ToString()} is always big-endian and RFC 9562 compliant.

\subsection{Limitations}\label{limitations}

\begin{enumerate}
\def\labelenumi{\arabic{enumi}.}
\item
  \textbf{Deployment-scoped uniqueness:} SKIDs are unique within a
  deployment (same epoch, same key-ring) but not globally unique across
  independent deployments. Applications requiring cross-deployment
  identity merging must include deployment identifiers in their routing
  logic.
\item
  \textbf{Topology limits and coordination requirement:} The default
  profile supports at most 128 applications times 64 instances = 8,192
  distinct generators. Topology assignment (AppId and AppInstanceId)
  requires deployment-time coordination, unlike UUID V4/V7/CUID2 which
  achieve uniqueness without coordination. The rationale for this
  specific field sizing is discussed in Design Trade-offs. Reference
  implementation aims to provide a coordination application called
  \texttt{Nexus} for this purpose.
\item
  \textbf{Key dependency:} SKEIDs and Secure SKEIDs require key material
  for generation and parsing. Loss of key material makes existing SKEIDs
  unparseable, though the underlying 64-bit SKIDs in the database remain
  valid and recoverable.
\item
  \textbf{Security analysis scope:} The security analysis in this paper
  uses STRIDE-based threat modeling and quantitative probability
  analysis within the defense-in-depth architecture. It does not provide
  formal cryptographic reductions (e.g., proving that the composition of
  AES-PRP encryption, BLAKE3 MAC, and the collision guard mechanism
  achieves the intended security goals under standard assumptions). The
  Secure SKEID scheme inherits PRP security from the AES block cipher,
  and the uniqueness of SKEID plaintexts prevents ciphertext equality
  leakage. A formal compositional security proof is left as future work.
  The reference implementation provides test vectors in its test suites
  enabling independent verification.
\item
  \textbf{Epoch-scoped storage efficiency:} The 8-byte SKID storage
  efficiency claim applies within a single epoch (\textasciitilde68
  years). The 64-bit SKID encodes a 32-bit timestamp relative to the
  configured epoch but does not carry the epoch index itself. A database
  spanning multiple epochs has the following options.

  \begin{enumerate}
  \def\labelenumii{\arabic{enumii}.}
  \item
    A composite key storing the epoch byte separately alongside the
    SKID. This adds per-row overhead that negates the 8-byte claim.
  \item
    Migration to the 16-byte SKEID, which includes the epoch at byte 0.
    This also negates the 8-byte claim.
  \item
    Epoch-partitioned storage, where the table or archive name
    implicitly carries the epoch context (e.g., \texttt{Entity\_Epoch0},
    \texttt{Entity\_Epoch1}). This preserves 8-byte rows by encoding the
    epoch in the partition structure rather than in each record.
  \end{enumerate}
\end{enumerate}

Since 68 years exceeds the operational lifespan of most modern software
systems, this constraint does not diminish practical utility.

\begin{enumerate}
\def\labelenumi{\arabic{enumi}.}
\setcounter{enumi}{5}
\tightlist
\item
  \textbf{Secure SKEID UUID format compliance:} A Secure SKEID is a
  valid UUID in length and string format (32 hexadecimal digits in
  \texttt{8-4-4-4-12} grouping) and is accepted by standard data-type
  parsers (e.g., PostgreSQL \texttt{uuid} columns, .NET \texttt{Guid},
  Java \texttt{UUID}). However, AES-256 encryption replaces the
  plaintext Version and Variant marker bytes with pseudorandom
  ciphertext. The encrypted output will almost never contain the
  \texttt{0x8} version nibble and \texttt{10xx} variant bits required by
  RFC 9562 Section 4.1 and Section 5.8. Strict RFC 9562 validators will
  reject a Secure SKEID as having an unknown version. Systems requiring
  strict RFC 9562 compliance at the validator level should use plaintext
  SKEIDs or accept Secure SKEIDs as opaque 128-bit values.
\end{enumerate}

\subsection{Future Work}\label{future-work}

Several extensions are planned or under consideration.

\begin{enumerate}
\def\labelenumi{\arabic{enumi}.}
\item
  \textbf{Epoch transition:} The mechanism for transitioning from epoch
  \texttt{0x00} to epoch \texttt{0x01} is not specified. Given that
  epoch \texttt{0x00} spans until approximately 2093 CE, this is left as
  future work.
\item
  \textbf{Key-ring rotation with multiple active key versions and
  graceful rollover:} The reference implementation currently uses a
  single key pair. The key-ring model and fallback algorithm are left as
  future work.
\item
  \textbf{Domain-tailored bit layouts:} Custom SKID profiles with
  different field widths (e.g., nanosecond or picosecond timestamp
  precision, single-year epoch windows for mission-scoped campaigns)
  would enable research workflows with different trade-off requirements.
  Configurable custom profiles are not standardized yet but under
  consideration.
\item
  \textbf{IETF standardization:} A draft specification prepared in
  Internet-Draft format (Kılıç, 2026) has been authored and is available
  for independent review, containing complete algorithms, CDDL
  definitions, and reproducible test vectors. Formal submission to the
  IETF would enable interoperability testing and community review.
\item
  \textbf{Formal compositional security proof:} A formal cryptographic
  reduction proving that the composition of single-block AES-PRP
  encryption, BLAKE3 MAC, and the collision guard mechanism achieves the
  intended security goals under standard assumptions would strengthen
  the security analysis beyond the STRIDE-based threat modeling and
  quantitative probability analysis provided in this paper.
\item
  \textbf{Cross-platform benchmarks:} The current benchmarks were
  conducted on ARM64 (Apple M2) with .NET 10. Benchmarks on x86-64
  architectures and other language runtimes would validate portability
  of the performance characteristics.
\item
  \textbf{Multi-language implementations:} The SKID design is
  language-agnostic. Implementations in Javascript, Java, Go, Rust,
  Python and other languages would expand the ecosystem and validate the
  specification's portability. These implementations are left to the
  open-source community.
\end{enumerate}

\section{Conclusions}\label{conclusions}

This paper presented Source Known Identifiers (SKIDs), a three-tier
identity system that simultaneously satisfies the proposed desired
identifier properties. Based on our literature survey, no existing
identifier scheme combines all six properties within a unified system.

The three-tier architecture aligns SKID (64-bit, database), SKEID
(128-bit, trusted internal), and Secure SKEID (128-bit, external) with
the trust boundaries commonly found in distributed applications.
Deterministic bidirectional transformations enable direct conversion
between tiers. The defense-in-depth security architecture combines
multiple verification layers to provide security far exceeding any
individual mechanism.

Benchmark results show that SKID and SKEID achieve superior performance
compared to UUID generation in the reference implementation. Secure
SKEID remains competitive while encapsulating desired identifier
properties with AES-256 encryption at approximately 1.4 times the
generation cost of UUID V7.

The open-source reference implementation, published as NuGet packages
with integration and unit test coverage, is available for community
review and adoption. A draft specification prepared in Internet-Draft
format is planned as future work to enable cross-platform implementation
and interoperability testing.

\section{Acknowledgments}\label{acknowledgments}

This work received no external financial support. GitHub CodeQL and
SonarCloud provided static analysis and code quality scanning that
strengthened the security posture of the implementation. CodeRabbit
provided AI code review that contributed to the overall code quality.

\section{Data Availability}\label{data-availability}

The source code for DRN-Project is available at
https://github.com/duranserkan/DRN-Project and archived on Zenodo (DOI:
\href{https://doi.org/10.5281/zenodo.19240397}{10.5281/zenodo.19240397}).
The archived release (v0.9.1) is available at
https://github.com/duranserkan/DRN-Project/tree/v0.9.1. Raw benchmark
data (CSV) and rendered reports (HTML) are included as six supplementary
files covering the primary benchmark, saturation benchmark, and hash
comparison benchmark.

\section{Competing Interests}\label{competing-interests}

The author declares no competing interests.

\section{Funding}\label{funding}

This work received no external financial support.

\section{Author Contributions}\label{author-contributions}

Duran Serkan Kılıç conducted the research, conceived the design,
implemented the software, conducted the benchmarks, and wrote the
manuscript.

\newpage

\section*{References}\label{references}
\addcontentsline{toc}{section}{References}

\protect\phantomsection\label{refs}
\begin{CSLReferences}{1}{0}
\bibitem[\citeproctext]{ref-benchmarkdotnet}
{Akinshin A, others}. 2013.
\href{https://github.com/dotnet/BenchmarkDotNet}{{BenchmarkDotNet:
Powerful .NET library for benchmarking}}.

\bibitem[\citeproctext]{ref-bellare1998}
Bellare M, Krovetz T, Rogaway P. 1998. {Luby-Rackoff Backwards:
Increasing Security by Making Block Ciphers Non-Invertible}. In:
\emph{Advances in cryptology -- EUROCRYPT '98}. LNCS. 266--280. DOI:
\href{https://doi.org/10.1007/BFb0054132}{10.1007/BFb0054132}.

\bibitem[\citeproctext]{ref-bonnetain2019}
Bonnetain X, Naya-Plasencia M, Schrottenloher A. 2019. Quantum security
analysis of {AES}. \emph{IACR Transactions on Symmetric Cryptology}
2019:55--93. DOI:
\href{https://doi.org/10.13154/tosc.v2019.i2.55-93}{10.13154/tosc.v2019.i2.55-93}.

\bibitem[\citeproctext]{ref-nistsp800107}
Dang Q. 2012. \emph{{Recommendation for Applications Using Approved Hash
Algorithms}}. NIST. DOI:
\href{https://doi.org/10.6028/NIST.SP.800-107r1}{10.6028/NIST.SP.800-107r1}.

\bibitem[\citeproctext]{ref-rfc9562}
Davis K, Peabody B, Leach P. 2024. \emph{{Universally Unique IDentifiers
(UUIDs)}}. IETF. DOI:
\href{https://doi.org/10.17487/RFC9562}{10.17487/RFC9562}.

\bibitem[\citeproctext]{ref-cuid-deprecated}
Elliott E. 2012. \href{https://github.com/paralleldrive/cuid}{{CUID:
Collision-resistant ids (deprecated)}}.

\bibitem[\citeproctext]{ref-cuid2}
Elliott E. 2022. \href{https://github.com/paralleldrive/cuid2}{{CUID2:
Secure, collision-resistant ids}}.

\bibitem[\citeproctext]{ref-evans2003}
Evans E. 2003.
\emph{\href{https://www.oreilly.com/library/view/domain-driven-design-tackling/0321125215/}{{Domain-Driven
Design: Tackling Complexity in the Heart of Software}}}. Addison-Wesley.

\bibitem[\citeproctext]{ref-ulid}
Feerasta A. 2016. \href{https://github.com/ulid/spec}{{ULID: Universally
Unique Lexicographically Sortable Identifier}}.

\bibitem[\citeproctext]{ref-ieee1588}
IEEE. 2020. \emph{{IEEE Standard for a Precision Clock Synchronization
Protocol for Networked Measurement and Control Systems}}. DOI:
\href{https://doi.org/10.1109/IEEESTD.2020.9120376}{10.1109/IEEESTD.2020.9120376}.

\bibitem[\citeproctext]{ref-brooks1995}
Jr. FPB. 1995.
\emph{\href{https://www.oreilly.com/library/view/mythical-man-month-the/0201835959/}{{The
Mythical Man-Month: Essays on Software Engineering}}}. Addison-Wesley.

\bibitem[\citeproctext]{ref-drn-project}
Kılıç DS. 2023. {DRN-Project: Distributed Reliable .Net}. DOI:
\href{https://doi.org/10.5281/zenodo.19240397}{10.5281/zenodo.19240397}.

\bibitem[\citeproctext]{ref-draft-skid}
Kılıç DS. 2026. {Source Known Identifiers (SKIDs): Technical
Specification Draft}. DOI:
\href{https://doi.org/10.5281/zenodo.19240397}{10.5281/zenodo.19240397}.

\bibitem[\citeproctext]{ref-rfc5905}
Mills D, Martin J, Burbank J, Kasch W. 2010. \emph{{Network Time
Protocol Version 4: Protocol and Algorithms Specification}}. IETF. DOI:
\href{https://doi.org/10.17487/RFC5905}{10.17487/RFC5905}.

\bibitem[\citeproctext]{ref-cwe639}
MITRE Corporation. 2025a.
\href{https://cwe.mitre.org/data/definitions/639.html}{{CWE-639:
Authorization Bypass Through User-Controlled Key}}.

\bibitem[\citeproctext]{ref-cwe345}
MITRE Corporation. 2025b.
\href{https://cwe.mitre.org/data/definitions/345.html}{{CWE-345:
Insufficient Verification of Data Authenticity}}.

\bibitem[\citeproctext]{ref-nistir8319}
Mouha N. 2021. \emph{{Review of the Advanced Encryption Standard}}.
NIST. DOI:
\href{https://doi.org/10.6028/NIST.IR.8319}{10.6028/NIST.IR.8319}.

\bibitem[\citeproctext]{ref-fips197}
National Institute of Standards and Technology. 2001. \emph{{Advanced
Encryption Standard (AES)}}. NIST. DOI:
\href{https://doi.org/10.6028/NIST.FIPS.197-upd1}{10.6028/NIST.FIPS.197-upd1}.

\bibitem[\citeproctext]{ref-nistsp800107withdrawal}
NIST. 2022.
\href{https://csrc.nist.gov/news/2022/withdrawal-of-nist-sp-800-107-revision-1}{{Withdrawal
of NIST Special Publication 800-107 Revision 1}}.

\bibitem[\citeproctext]{ref-blake3}
O'Connor J, Aumasson J-P, Neves S, Wilcox-O'Hearn Z. 2020.
\href{https://github.com/BLAKE3-team/BLAKE3-specs/blob/master/blake3.pdf}{{BLAKE3:
One function, fast everywhere}}.

\bibitem[\citeproctext]{ref-ruggles1947}
Ruggles R, Brodie H. 1947. {An Empirical Approach to Economic
Intelligence in World War II}. \emph{Journal of the American Statistical
Association} 42:72--91. DOI:
\href{https://doi.org/10.1080/01621459.1947.10501915}{10.1080/01621459.1947.10501915}.

\bibitem[\citeproctext]{ref-ksuid}
Segment. 2017. \href{https://github.com/segmentio/ksuid}{{KSUID:
K-Sortable Globally Unique IDs}}.

\bibitem[\citeproctext]{ref-shostack2014}
Shostack A. 2014.
\emph{\href{https://www.wiley.com/en-us/Threat+Modeling\%3A+Designing+for+Security-p-9781118809990}{{Threat
Modeling: Designing for Security}}}. Wiley.

\bibitem[\citeproctext]{ref-nistsp800224}
Turan MS, Brandão LTAN. 2024. \emph{{Keyed-Hash Message Authentication
Code (HMAC): Specification of HMAC and Recommendations for Message
Authentication}}. NIST. DOI:
\href{https://doi.org/10.6028/NIST.SP.800-224.ipd}{10.6028/NIST.SP.800-224.ipd}.

\bibitem[\citeproctext]{ref-snowflake}
Twitter Engineering. 2010.
\href{https://github.com/twitter-archive/snowflake}{{Announcing
Snowflake}}.

\end{CSLReferences}

\end{document}
